\chapter{场景编辑与操控}
\label{ch:09}

\paragraph*{本章阅读锚点}
本章的锚点是可控修改。原始 3DGS 优化后通常是固定场景；编辑方法要把语义、局部几何和纹理变化绑定到可操作的高斯集合或隐式控制信号上。


3D Gaussian Splatting 将场景表示为数百万个不可见的椭球——你可以渲染它们，但无法选中、点击或拖拽它们。当一位设计师想要"把这栋建筑换成哥特式风格"或"移除画面中的行人"时，面对的不是一张可编辑的图层，而是一堆散落在三维空间中的高斯基元。这一根本矛盾催生了 3DGS 编辑领域的研究：如何在缺乏拓扑结构的显式表示上实现精确、可控、多视角一致的编辑操作？

围绕这一核心挑战，研究从多个方向展开。第~\ref{sec:09_1}~节讨论\textbf{语义与文本引导编辑}——利用语言模型和扩散模型（Diffusion Model）实现"说一句话就能改场景"的编辑范式，包括基于得分蒸馏采样（SDS, Score Distillation Sampling）的迭代优化、迭代数据集更新（IDU, Iterative Dataset Update）策略、渐进式编辑和注意力定位等技术路线，以及交互式物体操控方法。第~\ref{sec:09_2}~节讨论\textbf{风格化与场景修复}——从优化式风格迁移到零样本前馈风格化，再到多视角一致的物体移除和空洞填充。第~\ref{sec:09_3}~节讨论\textbf{几何变形与其他方法}——笼式变形、网格代理、拖拽式操控等几何编辑技术，以及动态场景编辑和跨模态应用等新兴方向。

\begin{figure}[htbp]
\centering
\resizebox{0.95\textwidth}{!}{\begin{tikzpicture}[
  root/.style={rectangle, rounded corners=6pt, draw=techblue, fill=techbluebg, line width=1.2pt, text width=3.8cm, minimum height=0.9cm, align=center, font=\small\bfseries},
  method/.style={rectangle, rounded corners=3pt, draw=techgray!60, fill=white, text width=2.8cm, minimum height=0.55cm, align=left, font=\scriptsize, inner sep=3pt},
  edge/.style={draw=techgray!70, line width=0.6pt, -{stealth[length=4pt]}},
]
\node[root] (root) {3DGS 场景编辑与操控};

% Category nodes with explicit colors
\node[rectangle, rounded corners=4pt, draw=techblue, fill=techbluebg, line width=0.8pt, text width=2.8cm, minimum height=0.7cm, align=center, font=\footnotesize\bfseries, below left=1.2cm and 6.4cm of root] (sem) {语义/文本引导};
\node[rectangle, rounded corners=4pt, draw=techgreen, fill=techgreenbg, line width=0.8pt, text width=2.8cm, minimum height=0.7cm, align=center, font=\footnotesize\bfseries, below left=1.2cm and 3.2cm of root] (inter) {交互式操控};
\node[rectangle, rounded corners=4pt, draw=techorange, fill=techorangebg, line width=0.8pt, text width=2.8cm, minimum height=0.7cm, align=center, font=\footnotesize\bfseries, below left=1.2cm and 0cm of root] (style) {风格化};
\node[rectangle, rounded corners=4pt, draw=techpurple, fill=techpurplebg, line width=0.8pt, text width=2.8cm, minimum height=0.7cm, align=center, font=\footnotesize\bfseries, below right=1.2cm and 3.2cm of root] (inpaint) {场景修复};
\node[rectangle, rounded corners=4pt, draw=techgray, fill=techbg, line width=0.8pt, text width=2.8cm, minimum height=0.7cm, align=center, font=\footnotesize\bfseries, below right=1.2cm and 6.4cm of root] (deform) {几何变形};

\foreach \n in {sem,inter,style,inpaint,deform} {
  \draw[edge] (root.south) -- (\n.north);
}

% Semantic sub-methods
\node[method, below=0.3cm of sem.south west, anchor=north west] (s1) {\textbf{SDS 优化}: GSEdit, DreamCatalyst};
\node[method, below=0.08cm of s1] (s2) {\textbf{IDU}: GaussianEditor, DGE, TIGER};
\node[method, below=0.08cm of s2] (s3) {\textbf{渐进定位}: ProEdit, GSEditPro};
\node[method, below=0.08cm of s3] (s4) {\textbf{架构创新}: VF-Editor, DisCo3D};

% Interactive sub-methods
\node[method, below=0.3cm of inter.south west, anchor=north west] (i1) {\textbf{语义解耦}: 3DitScene, GaussianBlock};
\node[method, below=0.08cm of i1] (i2) {\textbf{分割操控}: Point'n Move, iSegMan};
\node[method, below=0.08cm of i2] (i3) {\textbf{颜色编辑}: ReCoGS, VIRGi};

% Stylization sub-methods
\node[method, below=0.3cm of style.south west, anchor=north west] (st1) {\textbf{优化式}: StylizedGS, StyleSplat};
\node[method, below=0.08cm of st1] (st2) {\textbf{分布匹配}: WaSt-3D};
\node[method, below=0.08cm of st2] (st3) {\textbf{前馈/零样本}: FPGS, AnyStyle};
\node[method, below=0.08cm of st3] (st4) {\textbf{多风格}: CLIPGaussian};

% Inpainting sub-methods
\node[method, below=0.3cm of inpaint.south west, anchor=north west] (p1) {\textbf{特征传播}: GScream, 3DGIC};
\node[method, below=0.08cm of p1] (p2) {\textbf{多视角扩散}: MVInpainter};
\node[method, below=0.08cm of p2] (p3) {\textbf{效率}: HOMER, GOR-IS};

% Deformation sub-methods
\node[method, below=0.3cm of deform.south west, anchor=north west] (d1) {\textbf{笼式}: GSDeformer, CAGE-GS};
\node[method, below=0.08cm of d1] (d2) {\textbf{网格代理}: Mani-GS, D-MiSo};
\node[method, below=0.08cm of d2] (d3) {\textbf{拖拽式}: DYG, ARAP-GS};

\end{tikzpicture}
}
\caption{3DGS 场景编辑方法的六大方向及其代表性工作。语义引导编辑形成四条技术路线（SDS优化、IDU、渐进定位、架构创新），交互式操控追求直接控制体验，风格化从优化式演进到前馈式，场景修复关注多视角一致性，几何变形通过笼式、网格代理和拖拽式三种策略实现。}
\label{fig:ch09_taxonomy}
\end{figure}

%=====================================================================
\section{语义与文本引导编辑}
\label{sec:09_1}

本节讨论利用语言模型和扩散模型实现"说一句话就能改场景"的编辑范式，涵盖文本/语义引导编辑（第~\ref{sec:09_1}~节）和交互式物体操控两个子方向。

\subsection{语义驱动的3DGS编辑}
\label{ss:09_1}

本节讨论两条技术路线：（1）文本/语义引导编辑；（2）交互式物体操控。

\subsubsection{文本/语义引导编辑}

本路线介绍利用文本指令或语义信息驱动 3DGS 场景编辑的方法。3DGS 场景缺乏语义结构：每个高斯基元只存储位置、颜色和不透明度，没有"这是桌子"或"那是窗户"的概念。当用户发出"把桌子变成木头纹理"的指令时，系统需要同时解决两个子问题：定位"桌子"对应哪些高斯基元；将"木头纹理"这一编辑意图转化为多视角一致的外观变化。

早期方法主要沿用 NeRF 编辑的思路，利用 2D 扩散模型作为编辑先验，通过得分蒸馏采样将 2D 编辑能力"注入"3D 表示。然而，SDS 的随机性导致不同视角产生不同编辑结果，引发多视角不一致问题。后续工作从不同角度回应了这一核心矛盾，形成了四条主要技术路线。

\paragraph{SDS 迭代优化路线}

最早的文本编辑方法直接将 SDS 损失应用于 3DGS 参数优化。

\textbf{动机。}3DGS 缺乏神经网络的记忆与缓冲能力，在面对扩散模型的随机梯度时容易产生不稳定更新，直接套用 NeRF 的 SDS 编辑范式效果有限。以 DreamFusion 的 SDS 为基线，其核心思想是：在每个训练步随机采样相机 $p$ 渲染图像 $I^p_{\text{RGB}}$，对渲染结果加噪后用扩散模型的 U-Net 预测噪声残差，再将残差梯度反传至 3DGS 参数。GSEdit（2024.03）~\cite{2403_05154}将此框架适配到 3DGS 编辑场景：采用 InstructPix2Pix（IP2P）作为编辑扩散模型——IP2P 是一种条件扩散模型，同时接受原始图像和文本指令作为输入，能在保持图像整体布局的同时按指令修改局部内容。在每个训练步，GSEdit 随机采样相机渲染当前场景，将渲染图像与编辑文本提示一同输入 IP2P。IP2P 的噪声预测与实际添加噪声之间的残差被用来引导 3DGS 参数的更新方向。直觉上，这个残差告诉优化器"当前渲染结果距离'按文本编辑后的样子'还差多少"。

其 SDS 损失定义为
\begin{equation}
  \mathcal{L}_{\text{SDS}} = \mathbb{E}_{t \sim \mathcal{U}(1,T)} \left[ w(t) \left( \hat{\epsilon}_\phi(z_t; t, c_I, c_T) - \epsilon \right) \right],
\end{equation}

其中 $z_t$ 是对渲染图像添加噪声 $\epsilon$ 后的隐编码（latent code），$\hat{\epsilon}_\phi$ 是 IP2P U-Net 在给定原始图像条件 $c_I$ 和编辑文本 $c_T$ 下预测的噪声，$w(t)$ 是时间步权重。残差项 $\hat{\epsilon}_\phi - \epsilon$ 编码了 IP2P 认为图像需要被编辑的程度与实际噪声之间的差异，该差异通过链式法则反传至 3DGS 的位置、颜色、透明度等参数。

GSEdit 的一个关键设计是时间步截断策略：将 $t$ 限制在 $[t_{\min}, t_{\max}]$ 区间内，并随训练推进线性降低该区间（参考 DreamTime 的策略）。之所以需要截断，是因为全时间步采样中高 $t$ 值对应的大噪声会引入过度随机扰动，容易破坏已有的几何结构。截断后，初始阶段加噪量较大（典型设置 $t_{\min}=0.02$, $t_{\max}=0.98$），允许模型在粗粒度上调整整体结构和外观。随着训练推进，$t_{\max}$ 逐步降低，加噪量减小，优化焦点转向细粒度纹理精修。整个编辑过程仅需约 500 步迭代，在消费级 GPU 上几分钟内即可完成，比同期 NeRF 编辑方法（如 Instruct-NeRF2NeRF 的约一小时）快一个数量级。然而，标准 SDS 框架存在两个根本缺陷：其一，省略了 U-Net Jacobian 项以保证计算效率，导致梯度估计偏差；其二，不同视角的 SDS 采样独立进行、时间步随机选取，缺乏跨视角协调，导致编辑结果出现纹理模糊和细节丢失。

与 GSEdit 直接使用 SDS 残差作为编辑梯度不同，DreamCatalyst（2024.07）~\cite{2407_11394}从扩散采样理论出发，指出了上述两个缺陷的根源——标准 SDS 将扩散模型当作简单的分数函数估计器，忽略了扩散反向过程中时间步的语义角色差异。

\textbf{动机。}GSEdit 的两个根本缺陷（省略 U-Net Jacobian 导致梯度偏差、跨视角时间步独立采样导致不一致）源于同一根因：SDS 忽略了扩散反向过程中时间步的语义差异；DreamCatalyst 从采样理论出发将 SDS 重新解释为单步 DDIM-SDEdit，为权重函数提供理论指导。其核心观察是：Delta Denoising Score（DDS）的目标函数等价于单步 DDIM-SDEdit 采样，而沿递减时间步优化 DDS 目标恰好对应反向 SDEdit 过程。在此基础上，DreamCatalyst 提出后验蒸馏采样（PDS, Posterior Distillation Sampling）的一般形式：

\begin{equation}
  \mathcal{L}_{\text{PDS}} = \mathbb{E}_{t,\epsilon} \left[ \lambda(t) \, \mathcal{L}_{\text{DDS}} + \mu(t) \, R_{\text{Iden}} \right], \quad R_{\text{Iden}} = \left\| \hat{x}_0^{\text{tgt}} - x_0^{\text{src}} \right\|_2^2,
\end{equation}

其中 $\mathcal{L}_{\text{DDS}}$ 负责编辑方向，$R_{\text{Iden}}$ 通过约束目标渲染 $\hat{x}_0^{\text{tgt}}$ 与源图像 $x_0^{\text{src}}$ 的 $\ell_2$ 距离来保持身份一致性，$\lambda(t)$ 和 $\mu(t)$ 是时间步相关的权重函数。简而言之，DDS 项告诉模型"往哪里改"，身份保持项告诉模型"什么不能丢"。DreamCatalyst 的关键创新在于为这两个权重设计了符合扩散采样动态的专门形式：

\begin{equation}
  \lambda^*(t) = \beta \, e^{-t/T}, \quad \mu^*(t) = \alpha + \gamma \, t/T,
\end{equation}

其中 $\alpha=0.075$, $\beta=0.2$, $\gamma=0.8$（全书实验统一设置）。这一设计满足两个条件：大时间步时身份保持权重大（$\mu^*(T) = 0.875$），避免高噪声扰动下源特征信息丢失；小时间步时身份保持权重小（$\mu^*(0) = 0.075$），允许扩散模型自由合成精细细节。配合递减时间步采样策略（$t$ 从 $T$ 均匀递减至 $1$），整个优化过程对齐了扩散模型的反向去噪轨迹，使 3DGS 参数更新方向与 2D 扩散先验保持一致。该方法提供快速模式（约 23 倍加速）和高质量模式（约 8 倍加速），同时支持 NeRF 和 3DGS 编辑。这种从"分数估计"到"采样过程近似"的范式转换，为后续方法提供了重要的理论视角。

\paragraph{GaussianEditor: 语义追踪与层次化泼溅}

与 SDS 直接优化 3D 参数不同，迭代数据集更新（IDU, Iterative Dataset Update）策略将编辑分解为两个独立步骤：先用 2D 扩散模型编辑选定视角的图像，再用编辑后的图像作为新训练数据更新 3DGS。这一解耦设计让 2D 编辑模型和 3D 重建各自在擅长的领域工作。

GaussianEditor（2023.11）~\cite{2311_14521}是首个系统性 3DGS 编辑框架。其设计动机源于一个核心观察：3DGS 作为点云式显式表示，在面对扩散模型的随机梯度时缺乏神经网络的记忆与缓冲能力，容易导致高斯点产生剧烈而不稳定的更新。为解决这一问题，GaussianEditor 提出两个关键技术。

其一是\textbf{高斯语义追踪}（Gaussian Semantic Tracing）：为每个高斯基元引入语义属性 $m_{ij}$（第 $i$ 个高斯对第 $j$ 个语义标签的隶属度），通过逆渲染将 2D 分割掩码反投影至 3D 高斯空间，权重计算为 $w_i^j = o_i(p) \cdot T_i^j(p) \cdot M^j(p)$，其中 $o_i$ 为不透明度、$T_i$ 为透射率、$M^j$ 为像素语义掩码。在自适应致密化过程中，新分裂的高斯继承父节点的语义标签，保证语义掩码在整个训练过程中持续有效——这克服了静态 2D/3D 掩码随场景演化而失效的问题。语义追踪还带来一个重要优势：在复杂场景中，2D 扩散模型往往难以有效编辑小物体；借助语义追踪，可以仅渲染目标物体并输入扩散模型，从而获得更精确的编辑监督。

其二是\textbf{层次化高斯泼溅}（HGS, Hierarchical Gaussian Splatting）：按高斯的致密化轮次划分代际，初始高斯为第 0 代，第 $k$ 轮致密化产生的新高斯标记为第 $k$ 代。每代高斯施加不同强度的锚点损失以控制自由度：

\begin{equation}
  \mathcal{L}^P_{\text{anchor}} = \sum_{i=0}^{n} \lambda_i (P_i - \hat{P}_i)^2, \quad \mathcal{L} = \mathcal{L}_{\text{Edit}} + \sum_{P \in \{x, s, q, \alpha, c\}} \beta_P \, \mathcal{L}^P_{\text{anchor}},
\end{equation}

其中 $P$ 为高斯的某一属性（位置 $x$、缩放 $s$、旋转四元数 $q$、不透明度 $\alpha$、颜色 $c$），$\hat{P}$ 为锚点状态，$\lambda_i$ 随代际递增——老一代高斯的锚点约束更强、自由度更低。直觉上，这一机制使早期代际逐渐"凝固"为稳定几何，而最新致密化的高斯在几乎无约束状态下雕琢细节，有效抑制了随机梯度导致的发散。锚点并非只在初始化时记录一次，而是在每轮致密化时更新为当前高斯状态，使约束始终相对于最新的几何结构生效。此外，HGS 采用百分比致密化策略（每步仅对位置梯度排名前 $k\%$ 的高斯进行分裂），比手动设定阈值更鲁棒。

GaussianEditor 通过语义追踪和 HGS 来约束 3DGS 在迭代中的更新行为，但本质上仍属于 IDU 范式——每轮先用 IP2P 独立编辑单张视角图像，再以编辑结果更新 3DGS。由于 IP2P 是单目编辑器，从同一场景的不同视角独立采样编辑结果时，获得一致编辑的概率几乎为零。GaussianEditor 的 HGS 锚点损失和语义追踪只是在 3D 空间缓解这一不一致性，并未从根本上解决多视角编辑的统计一致性问题。

\paragraph{DGE: 图像空间多视角一致编辑}

DGE（Direct Gaussian Editor）（2024.04）~\cite{2404_18929}正是针对这一瓶颈提出了替代思路：与其在 3D 空间反复修补不一致的编辑结果，不如直接在图像空间保证多视角一致性，然后仅需将 3DGS 拟合到编辑后的图像即可完成 3D 编辑。与 GaussianEditor 通过 HGS 逐代凝固高斯属性来抑制迭代发散不同，DGE 将核心难题前移至 2D 编辑阶段——借鉴视频编辑中无需额外训练即可扩展图像生成器的思路，将多视角序列视为轨道视频，使多视角一致性等价于视频编辑中的时序一致性。DGE 将流程分为两步：\textbf{关键视角编辑}和\textbf{特征注入}。关键视角编辑阶段，在 IP2P 的自注意力层中引入时空注意力，使 $T$ 个视角的 query 同时关注所有视角的 key-value：

\begin{equation}
  \text{Attn}(\mathbf{Q}, \mathbf{K}, t) = \text{softmax}\!\left( \frac{\mathbf{Q}_t [\mathbf{K}_1, \ldots, \mathbf{K}_T]^\top}{\sqrt{d}} \right),
\end{equation}

但该操作计算开销随视角数二次增长，因此实际执行中仅对随机采样的关键视角子集 $\mathcal{K} \subset \{1, \ldots, T\}$ 执行。在特征注入阶段，DGE 利用\textbf{极线约束}（Epipolar Constraint）将关键视角的编辑特征传播到其余视角：给定非关键视角 $t'$ 的特征图 $\mathbf{F}_{t'}$ 和最近关键视角 $k^*$ 的特征图 $\mathbf{F}_{k^*}$，对应关系通过基础矩阵 $F$ 约束搜索空间获得：

\begin{equation}
  \mathbf{M}_{t'}[u] = \arg\min_{v,\; Fv \parallel u} D\!\left( \mathbf{F}_{t'}[u],\, \mathbf{F}_{k^*}[v] \right),
\end{equation}

其中 $D$ 为余弦距离，$Fv \parallel u$ 表示点 $v$ 在关键视角中的对应点必须位于极线 $Fu$ 上。这一约束将对应搜索从二维平面压缩到一维极线上，降低了误匹配率。最终 DGE 通过渲染损失 $\theta' = \arg\min_\theta \sum_{t=1}^{T} \| I'_t - R(\theta, \pi_t) \|$ 直接优化 3DGS 参数，整个编辑流程无需额外训练扩散模型。

\paragraph{TIGER: 一致性得分蒸馏}

上述方法（GaussianEditor 和 DGE）均遵循 IDU 范式——先独立编辑 2D 视角图像、再更新 3DGS，但由于每帧编辑相互独立，不可避免地导致过平滑和多面 Janus 问题。TIGER（2024.05）~\cite{2405_14455}从这一瓶颈出发，提出了另一种解决思路：不再逐帧独立编辑后拼合，而是通过得分蒸馏（SDS）同时驱动所有视角的更新，并引入多视角扩散模型从源头保证几何一致性，从而绕过 IDU 范式中 2D 不一致向 3D 传播的根本难题。

\textbf{动机。}IDU 范式中每帧独立编辑导致 2D 不一致向 3D 传播，即使 GaussianEditor 的 HGS 锚点和 DGE 的极线约束也只是在空间层面修补不一致性，而非从源头消除；TIGER 通过同时驱动所有视角的更新，从根本上避免不一致的产生。

与 GaussianEditor 依赖 2D 分割掩码反投影为高斯赋予语义标签（自顶向下）不同，TIGER 为每个高斯嵌入 64 维语言向量 $L \in \mathbb{R}^{64}$，采用\textbf{自底向上}的语言特征提取策略：先用 MaskCLIP 生成低分辨率 patch 级语言特征（包含跨语义边界的全局上下文），再用 FeatUp 上采样至像素级，经 PCA（Principal Component Analysis，主成分分析）降维至 64 维后作为高斯语言嵌入的监督信号。为进一步精修语义边界的不均匀区域，TIGER 引入 SAM（Segment Anything Model）在最细层级生成二值掩码，对掩码内语言特征取加权平均，使边界更加锐利——这与 LangSplat/Gaussian Grouping 的自顶向下策略（先分割再在掩码内提取 CLIP 特征）有本质区别，后者因掩码内缺乏外部上下文而限制开放词汇查询能力。

在编辑阶段，TIGER 提出\textbf{一致性得分蒸馏}（CSD），同时利用 2D 图像编辑扩散模型（InstructPix2Pix）和 4 视角多视角扩散模型（MVDream）的得分：

\begin{equation}
  \nabla_S \mathcal{L}_{\text{CSD}}(S) = \mathbb{E}_{v, i} \left[ \omega_1 \nabla_S \mathcal{L}_{\text{ip}} + \omega_2 \nabla_S \mathcal{L}_{\text{mv}} \right],
\end{equation}

其中 $\mathcal{L}_{\text{ip}}$ 和 $\mathcal{L}_{\text{mv}}$ 分别为 InstructPix2Pix 和 MVDream 在渲染视角上的噪声残差，$\omega_1$、$\omega_2$ 为动态权重。训练初期，$\omega_2$ 权重较高，优先由 MVDream 建立一致的几何结构；后期逐步增大 $\omega_1$，由 InstructPix2Pix 雕刻编辑细节。此外，TIGER 始终以原始 3DGS 渲染图像作为 InstructPix2Pix 的条件输入，避免高斯过度流动导致的不稳定更新；在噪声注入上采用极低噪声水平（$t \in [0.02, 0.2]$），确保每次更新的增量微小而稳定。每个高斯的更新幅度还受其与文本查询的相关度分数 $s_i = T^{*\top} G_i(L)$ 调制，相关度高的高斯更新快、无关高斯几乎不更新，实现精确的对象级编辑。

\paragraph{TrAME 与 3DEgo: IDU 理论与工程改进}

除了上述核心管线改进，还有两类工作从不同侧面推进了 IDU 方法。在\textbf{理论层面}，TrAME（2024.07，Trajectory-Anchored Multi-view Editing）~\cite{2407_02034}揭示了优化式和重建式 3D 编辑方法的等价关系，并据此设计轨迹锚定方案（TAS）和免调参的视角一致注意力控制模块——这一发现暗示 SDS 和 IDU 路线的性能差距可能源于实现细节而非根本范式差异。在\textbf{工程效率层面}，3DEgo（2024.07）~\cite{2407_10102}将传统三阶段流程简化为单阶段工作流，通过噪声混合器模块实现免训练的多视角一致 2D 编辑。基于注意力特征变形的方法（2024.12）~\cite{2412_07984}则走得更远——它仅使用单视角编辑作为主操作源，利用 3DGS 深度估计进行几何引导的注意力变形，将编辑一致传播到所有视角，无需同时处理多帧，大幅降低了计算开销。

\begin{figure}[htbp]
\centering
\resizebox{0.9\textwidth}{!}{\begin{tikzpicture}[
  box/.style={rectangle, rounded corners=3pt, draw=#1, fill=#1!8, line width=0.7pt, text width=2.6cm, minimum height=0.6cm, align=center, font=\scriptsize\bfseries},
  arr/.style={-{stealth[length=3.5pt]}, line width=0.5pt, draw=techgray},
  lbl/.style={font=\scriptsize, text=techgray},
]
% SDS pipeline (top)
\node[font=\small\bfseries, text=techblue] at (-5.5, 3.8) {SDS 迭代优化};

\node[box=techblue] (s1) at (-5.5, 2.5) {3DGS 渲染\\$I^p = R(\theta, p)$};
\node[box=techblue] (s2) at (-5.5, 1.5) {加噪 $z_t$\\扩散模型\\预测残差};
\node[box=techblue] (s3) at (-5.5, 0.5) {SDS 梯度\\$\nabla_\theta \mathcal{L}_{\text{SDS}}$};
\node[box=techorange] (s4) at (-5.5, -0.5) {更新 3DGS\\参数 $\theta$};

\draw[arr] (s1) -- (s2);
\draw[arr] (s2) -- (s3);
\draw[arr] (s3) -- (s4);
\draw[arr, dashed, bend left=60] (s4.east) to node[lbl, right] {迭代循环} (s1.east);

% IDU pipeline (bottom)
\node[font=\small\bfseries, text=techgreen] at (2, 3.8) {迭代数据集更新 (IDU)};

\node[box=techgreen] (d1) at (2, 2.5) {3DGS 渲染\\选定视角图像};
\node[box=techgreen] (d2) at (2, 1.5) {2D 扩散模型\\编辑图像\\(IP2P/\allowbreak ControlNet)};
\node[box=techgreen] (d3) at (2, 0.5) {编辑后图像\\作为新训练集};
\node[box=techorange] (d4) at (2, -0.5) {3DGS 重建\\$\min_\theta \|I'_t - R(\theta)\|$};

\draw[arr] (d1) -- (d2);
\draw[arr] (d2) -- (d3);
\draw[arr] (d3) -- (d4);
\draw[arr, dashed, bend right=60] (d4.west) to node[lbl, left] {迭代更新数据集} (d1.west);

% Key difference
\draw[{stealth[length=4pt]}-{stealth[length=4pt]}, thick, draw=techpurple] (-2.5, 1.5) -- (-0.5, 1.5)
  node[midway, above, font=\scriptsize\bfseries, text=techpurple] {核心差异};
\node[font=\scriptsize, text=techpurple, align=center] at (-1.5, 0.8) {SDS: 直接优化3D参数\\IDU: 解耦2D编辑+3D重建};

% Pros/cons
\node[font=\scriptsize, align=left, text=techblue] at (-5.5, -1.6) {
  \textbf{优点}: 端到端优化\\
  \textbf{缺点}: 梯度偏差, 视角不一致
};
\node[font=\scriptsize, align=left, text=techgreen] at (2, -1.6) {
  \textbf{优点}: 各模块各司其职\\
  \textbf{缺点}: 2D不一致→3D抖动
};

\end{tikzpicture}
}
\caption{3DGS 文本引导编辑的两种核心范式对比。SDS 迭代优化（左）直接将扩散模型的梯度信号反传至 3DGS 参数，端到端但有梯度偏差；迭代数据集更新（右）将 2D 编辑和 3D 重建解耦为独立步骤，各模块各司其职但需处理 2D 不一致性传播。}
\label{fig:ch09_sds_vs_idu}
\end{figure}

\paragraph{ProEdit: 渐进式文本插值定位}

SDS 和 IDU 方法面临一个共同的精度问题：扩散模型的可行输出空间（FOS, Feasible Output Space）过大，不同视角可能产生截然不同的编辑结果。与前两种方法在优化层面寻求一致性不同，渐进式方法从任务层面切入——通过将编辑任务分解为多个子任务来控制每步的输出空间大小。

\textbf{动机。}扩散模型的 FOS 越大，不同视角的编辑结果越容易分歧；通过文本嵌入插值逐步缩小编辑跨度，每步的 FOS 足够小以确保多视角自然趋于一致。

ProEdit（2024.11）~\cite{2411_05006}是该路线的代表方法，其核心洞察是多视角不一致性源于扩散模型的大 FOS：一次性从猫变成狗的编辑空间太大，不同视角可能走向完全不同的狗。但若先变成猫狗混合体，再逐步过渡到狗，每步的编辑自由度就足够小，多视角自然趋于一致。具体而言，ProEdit 将编辑任务 $T(s, e=E(p))$ 定义为在原始场景 $s$ 上应用文本指令 $p$，其中 $e$ 为冻结文本编码器 $E(\cdot)$ 计算的文本嵌入。子任务通过文本嵌入插值来构造——例如 $r{=}0.5$ 对应的子任务可直觉理解为"让他半微笑"：

\begin{equation}
    S(s, r) = T\big(s,\; e(r)\big), \quad e(r) = r \cdot E(p) + (1-r) \cdot E(\emptyset),\; r \in [0,1]
\end{equation}

其中 $E(\emptyset)$ 为空提示的嵌入（对应恒等映射），$r$ 为插值比例。当 $r=0$ 时场景不变，$r=1$ 时完成完整编辑。这一连续子任务空间将大跨度编辑拆分为序列 $\{r_0, r_1, \ldots, r_n\}$（$r_0=0, r_n=1$），每步的 FOS 足够小以确保多视角一致。

\textbf{难度感知调度。}为使各子任务难度均衡，ProEdit 需要度量相邻子任务间的难度差异。FOS 的大小本身难以直接计算，因此 ProEdit 用 LPIPS（Learned Perceptual Image Patch Similarity，感知图像块相似度）感知距离来近似：先用 2D 扩散模型分别对原始图像施加 $e(r_i)$ 和 $e(r_j)$ 两种编辑，再用 LPIPS 度量两组编辑图像的差异 $d(r_i, r_j)$——差异越大意味着该区间编辑自由度越高。基于此度量，ProEdit 以递归二分法自适应确定子任务序列：在 $[r_a, r_b]$ 区间上检查 $d(r_a, r_b) \leq d_{\text{threshold}}$，若超过阈值则取中点 $r_m$ 递归划分 $[r_a, r_m]$ 和 $[r_m, r_b]$，直到所有相邻子任务满足阈值约束，保证每步编辑难度均衡。

\textbf{3DGS 自适应训练。}由于 3DGS 原为多视角一致图像的重建设计，直接在不一致编辑图像上训练会导致高斯数量持续膨胀直至内存溢出。ProEdit 提出自适应高斯创建策略来解决这一问题。在每个子任务开始时，将所有高斯的透明度重置为剔除阈值，通过若干轮训练隐式识别位于物体表面的高斯（这些高斯的透明度在训练中重新升高而得以保留），而几何不正确的高斯透明度保持低位并在下次维护阶段被剔除。同时，根据现有高斯数量与已剔除数量动态调节梯度阈值，控制每轮新增高斯的数量，避免不一致性导致的高斯数量爆炸。此外，鉴于扩散模型推理（约 2 秒/张）与 3DGS 训练（约 0.02 秒/轮）速度差异悬殊，整个流程采用双 GPU 解耦调度：GPU1 负责扩散模型推理生成编辑图像，GPU2 负责 3DGS 训练，二者以缓冲区异步通信。

\paragraph{GSEditPro: 语义标签空间定位}

ProEdit 通过文本插值缩小 FOS 来隐式缓解多视角不一致，但并未显式约束编辑的空间位置。当编辑目标仅涉及场景中的特定物体时，缺乏空间定位意味着 SDS 梯度会无差别地作用于所有高斯，导致编辑泄漏到无关区域。GSEditPro（2024.11）~\cite{2411_10033}正是从这一定位精度缺口出发，在 3D 高斯层面引入显式的语义标签，将编辑严格约束在目标区域。

\textbf{动机。}ProEdit 在 FOS 维度缩小编辑自由度，但缺乏空间维度的定位能力；GSEditPro 通过交叉注意力机制为每个高斯分配二值语义标签，在空间维度锁定编辑范围，二者互补。

GSEditPro 的定位管线由三个阶段递进构成：交叉注意力定位、SAM 精确分割、3D 反投影标记。首先，为增强扩散模型对目标对象的注意力准确性，GSEditPro 通过 DreamBooth 以目标编辑提示作为类别提示微调 Stable Diffusion——其保留损失促使扩散模型将该类别作为默认生成风格，从而提升注意力图的质量。在 Img2Img 生成过程中，提取目标词的交叉注意力图作为编辑区域的粗略定位。由于编辑目标分为添加新物体与修改已有物体两种情形，GSEditPro 根据目标物体是否已存在于场景中采用不同策略：(1) 添加新物体时，注意力图经阈值 0.5 二值化后使用 DBSCAN 聚类滤除离群点，直接作为 2D 掩码；(2) 编辑已存在物体时，选取注意力值最高的 5 个点作为正样本、最低的 3 个点作为负样本，输入 SAM 获得精确分割掩码——这比直接使用语言提示驱动 SAM 更可靠，因为 SAM 的语言分割在不同视角间不稳定。最后，这些 2D 掩码通过可微渲染反投影到 3D 空间，仅允许梯度在标记高斯上回传，在训练开始前即完成显式的二值语义标签分配。

编辑优化采用 SDS 损失 $\mathcal{L}_{\text{SDS}}$（由 DreamBooth 微调后的扩散模型计算）结合锚定损失 $\mathcal{L}_{\text{anchor}}^P$（约束高斯位置、缩放、旋转、透明度和颜色不偏离历史状态）以及像素级保持损失。由于高斯的强流动性，即使梯度仅在标记高斯上回传，渲染结果中非编辑区域仍可能因高斯漂移而改变。为此 GSEditPro 构造伪真值图像 $I_{\text{pgt}} = M \cdot I_{\text{edit}} + (1-M) \cdot I_{\text{ori}}$，由当前编辑结果与原始图像按掩码拼接而成，确保非编辑区域不变。值得注意的是，掩码 $M$ 本身在优化过程中分两阶段演进：前 2000 次迭代使用定位模块生成的静态掩码构造粗略伪真值；此后当编辑结果已初步成形时，重新调用定位模块获取动态掩码，使伪真值随训练自适应更新。像素级保持损失为：

\begin{equation}
    \mathcal{L}_{\text{Preservation}} = \lambda_{l1}\,\mathcal{L}_1(I_{\text{edit}}, I_{\text{pgt}}) + \lambda_{\text{ssim}}\,\mathcal{L}_{\text{D-SSIM}}(I_{\text{edit}}, I_{\text{pgt}})
\end{equation}

总损失为 $\mathcal{L} = \lambda_{\text{sds}}\mathcal{L}_{\text{SDS}} + \mathcal{L}_{\text{Preservation}} + \sum_{P \in \{x,s,q,\alpha,c\}} \lambda_P \mathcal{L}_{\text{anchor}}^P$，其中锚定损失的强度随迭代次数递增，逐步收紧对高斯属性的约束。

\paragraph{EditSplat: 多视角融合与注意力修剪}

ProEdit 通过渐进式子任务缩小 FOS 来隐式缓解多视角不一致——但 SDS 梯度在各视角独立回传、缺乏显式跨视角约束，不同视角编辑结果仍可能沿各自梯度方向发散；GSEditPro 通过语义标签在空间维度锁定编辑范围——但 2D 掩码反投影至 3D 后缺乏跨视角一致性校验，不同视角的标签可能冲突，且二值标签丢失连续重要性权重，优化缺乏精细控制。两条路线各自只解决了一个维度的问题——ProEdit 缺乏空间定位且多视角一致性是隐式的，GSEditPro 空间定位精确但跨视角一致性是隐式的。EditSplat（2024.12）~\cite{2412_11520}从这一双重缺口出发，同时从多视角一致性和高斯管理两个维度切入，提出多视角融合引导（MFG）和注意力引导修剪（AGT）两个互补机制。

\textbf{动机。}此前方法各自只解决一个维度（FOS 或空间定位），EditSplat 同时从两个维度入手：MFG 显式将多视角信息注入扩散过程，AGT 通过连续注意力权重和修剪策略管理高斯密度。

\textbf{多视角融合引导（MFG）。}MFG 首先用扩散模型编辑所有源图像，通过 ImageReward 评分剔除底部 15\% 低质量结果，选取与目标视角最接近的 top-5 相邻视角，利用 3DGS 深度图按深度值由远及近迭代 $\alpha$-混合，将多视角编辑结果投影融合到每个目标视角得到 $h_M$。对融合后稀疏的背景区域，MFG 使用 SAM 以源内容填充以实现无缝整合。融合图像 $h_M$ 通过 classifier-free guidance 直接注入扩散过程，无需额外训练、参数化层或重复扩散过程：

\begin{equation}
    \hat{\epsilon}_\theta(z_t, h_S, h_T, h_M) = \epsilon_\theta(\cdot,\emptyset,\emptyset) + s_T\,\Delta_T + s_M\,\Delta_M + s_S\,\Delta_S
\end{equation}

其中 $\Delta_T = [\epsilon_\theta(z_t, h_M, h_T) - \epsilon_\theta(z_t, h_M, \emptyset)]$，$\Delta_M$ 和 $\Delta_S$ 类似地捕获多视角融合图像与源图像的引导贡献，$s_T$、$s_M$、$s_S$ 为对应引导强度。通过平衡三种引导信号，MFG 引导扩散模型生成与多视角融合图像在纹理、颜色和几何上对齐的编辑结果，在单次扩散过程中即实现多视角一致性。

\textbf{注意力引导修剪（AGT）。}AGT 解决两个互补问题：(1) 预训练高斯过度保留源属性，阻碍编辑优化；(2) 此前方法使用 SAM 提取的 2D 掩码仅基于源图像而非编辑结果，可能遗漏语义相关区域。与 GSEditPro 使用交叉注意力生成二值标签（编辑/非编辑）不同，AGT 将多视角交叉注意力图通过逆渲染分配给每个 3D 高斯连续权重：

\begin{equation}
    w_j = \frac{1}{\sum_{v \in V}|S_{v,j}|} \sum_{v \in V} \sum_{s \in S_{v,j}} (Q_v K_v^T)_{s}
\end{equation}

其中 $S_{v,j}$ 为第 $j$ 个高斯在第 $v$ 个视角注意力图上影响的像素索引集合。通过跨视角聚合注意力权重，每个高斯获得反映整体编辑重要性的连续值，而非 GSEditPro 的二值标签。AGT 对高斯执行两级操作：低于阈值 $w_{\text{thres}}$ 的高斯梯度被置零以实现选择性优化，将更新集中于语义相关区域；同时，注意力权重最高的 top-$k\%$ 高斯被直接修剪——这似乎反直觉（删除最重要的编辑区域），但其动机是这些高斯对应需要大幅修改的区域，它们过度保留了源属性，在致密化过程中会生成更多携带源信息的高斯，因此直接修剪后重新致密化反而更高效。最终编辑损失在所有视角上最小化 $\mathcal{L}_{\text{edit}} = \text{LPIPS} + \mathcal{L}_1$（权重比 1:1）。

\paragraph{Mastering Regional 与互补一致性: 渐进式增量}

沿渐进式定位路线，后续方法从不同角度持续改进。Mastering Regional 方法（2025.07）~\cite{2507_05426}利用 2D 扩散编辑精确定位 3DGS 修改区域，结合深度先验初始化和序贯视角细化，在 SOTA 性能上实现 4 倍加速。互补与一致性信息方法（2025.03）~\cite{2503_11601}利用多视角互补信息增强编辑一致性。

\paragraph{ProGDF: out-of-loop 实时编辑}

近期工作从架构层面重新思考 3DGS 编辑范式。前三条路线（SDS 引导、迭代数据更新、渐进式优化）共享同一个范式瓶颈：编辑过程与优化过程交织在一起，即 in-the-loop 训练——每次编辑指令的修改都需要重新启动一轮完整优化（据 ProGDF 原文每次调整需数十分钟训练成本），用户无法实时交互。ProGDF（2024.12）~\cite{2412_08152}从根本上扭转了这一范式，将编辑从 in-the-loop 转变为 out-of-loop：训练一次、实时编辑。

\textbf{动机。}in-the-loop 训练的编辑范式使用户每次修改指令都需等待完整优化，无法实现交互式体验；ProGDF 用轻量级神经网络学会从原始场景到目标场景的整个编辑轨迹，用户只需拖动滑块选择轨迹上的时间点即可实时获得编辑结果。其直觉是：如果用一个轻量级神经网络学会从原始场景到目标场景的整个编辑轨迹，那么用户只需拖动滑块选择轨迹上的某个时间点，就能在 $0.02$ 秒内获得对应的编辑结果，无需再次优化。ProGDF 由两个组件协同实现这一目标：渐进式高斯泼溅（PGS, Progressive Gaussian Splatting）负责在优化轨迹上提取高质量中间结果；高斯微分场（GDF, Gaussian Differential Field）负责学习这些中间结果并参数化为滑块控制信号。

PGS 的核心挑战在于：常规优化轨迹的前期大量高斯朝不精确方向更新，导致纹理模糊和细节扭曲；后期虽然收敛但中间状态缺乏多样性，难以提供丰富的训练监督。为此，PGS 提出两项约束平滑优化轨迹。渐进约束 $L_{prog} = \beta^{t/s}\sum_{i=1}^{N}\|\delta_i^t\|$ 通过衰减系数 $\beta^{t/s}$（$0<\beta<1$，$s$ 用于缓解指数爆炸）自适应压制每步更新幅度，使场景从原始到目标的过渡更加平稳。渲染质量正则化 $L_{render} = -\|\nabla_{Laplacian}^2 R(\Theta^t,v)\|^2$ 利用拉普拉斯梯度衡量图像清晰度，梯度越大表示越清晰，以此显式约束中间结果的纹理质量。PGS 总损失为三者的加权和 $L_{PGS} = \lambda_1 L_{edit} + \lambda_2 L_{prog} + \lambda_3 L_{render}$，其中 $\lambda_1=1$、$\lambda_2=5$、$\lambda_3=1$，渐进约束的权重最高以充分平滑轨迹。优化过程中，PGS 将轨迹上的中间结果均匀采样后存入在线记忆库，供 GDF 训练时随机抽取。

GDF 是一个轻量级神经网络，以空间中任意点 $(x,y,z)$ 为输入，预测该位置高斯基元因编辑产生的属性偏移量 $(\delta\mu_i, \delta\Sigma_i, \delta\alpha_i, \delta c_i)$。为实现可控编辑，GDF 将优化轨迹的时间步离散化为 $k$ 个 bin，通过可学习变换 $H$ 将归一化控制信号 $\tau \in [0,1]$ 映射为偏移量，推理公式为 $(\delta\mu_i^t, \delta\Sigma_i^t, \delta\alpha_i^t, \delta c_i^t) = G(\Theta_i^0, H(\tau))$，其中 $\tau=0$ 对应零偏移（原始场景）、$\tau=1$ 对应完全编辑结果。编辑后的高斯通过逐元素叠加得到：$\Theta_i^{edit} = (\mu_i^0 + \delta\mu_i^t, \Sigma_i^0 + \delta\Sigma_i^t, \alpha_i^0 + \delta\alpha_i^t, c_i^0 + \delta c_i^t)$。为防止 GDF 在与 PGS 并行训练时发生灾难性遗忘，ProGDF 采用轨迹重采样策略——从记忆库中随机采样中间结果作为监督而非按序使用。GDF 损失函数为 $L_{GDF} = L_{pred} + L_{render}$，其中 $L_{pred}$ 衡量 GDF 预测编辑结果与对应中间结果的 $L_1$+LPIPS 差距。整个系统仅需一次 out-of-loop 训练，之后用户通过滑块界面在 $0.02$ 秒内实时输出可控编辑结果。

\paragraph{3DSceneEditor: 全 3D 管线交互式编辑}

沿架构创新路线，3DSceneEditor（2024.12）~\cite{2412_01583}是首个全 3D 管线的交互式场景编辑框架，直接在 3D 高斯基元层面操作，支持添加、移动、重新着色、替换和删除物体，集成预训练实例分割和 CLIP 零样本接地模块。

\textbf{动机。}前述方法多依赖 2D 编辑扩散模型间接驱动 3DGS 更新，编辑精度受限于 2D-3D 传播过程的不一致性；3DSceneEditor 直接在 3D 空间操作高斯基元，实现精确的物体级编辑。

\paragraph{VF-Editor: 前馈式原生 3DGS 编辑器}

直觉上，VF-Editor 就像一个"翻译官"——先学会理解 2D 编辑的"语言"，然后直接将编辑指令翻译为每个 3D 高斯的属性变化，无需反复试错。

\textbf{动机。}所有前述编辑方法均需逐场景优化，用户每次编辑都需等待分钟级的优化过程；VF-Editor 通过离线蒸馏多种 2D 编辑策略的知识到通用预测器，实现一次前向传播完成编辑。VF-Editor（2026.02）~\cite{2602_11638}提出首个前馈式原生 3DGS 编辑器，从根本上绕过了逐场景优化的瓶颈：其核心思想是训练一个通用的变化预测器 $P_\theta$，通过离线蒸馏 2D 编辑知识（IP2P、CtrlColor、DDPM inversion 等多种编辑策略），使推理时只需一次前向传播即可预测所有高斯基元的属性变化量，整个编辑过程约 $0.3$ 秒。但该设计存在另一层面的代价：离线蒸馏阶段需对每个编辑策略分别执行大规模知识迁移，训练通用预测器 $P_\theta$ 需消耗大量 GPU 显存与 epoch 数量，且蒸馏质量受限于教师模型的多视角一致性——若教师模型本身存在视角伪影，蒸馏后的预测器亦会继承该缺陷。$P_\theta$ 由三个组件构成：随机分词器 $T$、变化场生成模块 $M$ 和迭代并行解码函数 $F$。

\textbf{随机分词器 $T$}。3DGS 中不同场景的高斯数量差异巨大且空间分布极不均匀，无法直接送入固定输入尺寸的 transformer。$T$ 的解决方案是：从源场景 $X^s$ 中随机采样 $n$ 个高斯作为锚点，对每个锚点选取 $k-1$ 个空间最近邻高斯组成一组，将每组 $k$ 个高斯的属性拼接为维度 $k \cdot f$ 的 3D token（$f$ 为单个高斯基元的特征维度），再经 MLP 进行维度变换。选择随机采样而非点云处理中常用的最远点采样，是因为 3DGS 的稀疏边缘基元密度极低，最远点采样会过度选中外围稀疏点，而随机采样能得到更合理的锚点分布。

\textbf{变化场生成模块 $M$}。$M$ 的设计动机源于一个关键观察：2D 编辑方法固有的概率流（probabilistic flow）是多视角不一致的主要来源——不同视角使用不同随机噪声会导向不同的编辑结果。现有方法通常选择限制这种多样性以换取一致性，代价是编辑结果的丰富性降低。VF-Editor 采取了截然不同的策略：不在蒸馏时限制概率流，而是保留与概率流强相关的关键噪声 $z$（即 DDIM 推理的初始噪声），让 $P_\theta$ 学会存储概率流的可能输出。具体而言，$M$ 由堆叠的 transformer block 构成，将 3D token 与 $z$ 拼接后作为输入，CLIP 编码的编辑指令 $y$ 通过交叉注意力注入，生成变化场
\begin{equation}
f_\theta = M(T(X^s) \oplus z;\, y),\quad z \sim \mathcal{N}(0,I),
\end{equation}

其中 $\oplus$ 表示拼接操作。将 $z$ 纳入输入而非丢弃，使得单视角编辑结果的监督即可有效缓解多视角不一致——模型在充分训练后甚至能从已知视角的变化推断新视角的合理编辑。

\textbf{迭代并行解码函数 $F$}。$F$ 的任务是从全局变化场 $f_\theta$ 中高效提取每个高斯基元的属性变化量 $\Delta_i = (\delta\mu_i, \delta s_i, \delta\alpha_i, \delta c_i, \delta r_i)$。$F$ 采用无自注意力的 transformer 架构：以各高斯基元属性为 query、变化场为 key/value 进行交叉注意力，每个基元独立解码，复杂度保持 $O(N)$，可完全并行化。$F$ 故意不使用自注意力，因为全局信息聚合已由 $M$ 完成，$F$ 只需轻量级条件读取（conditional readout）。为缓解 3D 高斯倾向于改变外观属性而非移动几何位置的问题（源于渲染方程中位置与外观属性的强耦合 Jacobian），$F$ 将均值 $\delta\mu$ 与其余属性分开迭代：第一阶段 $F_1$ 先解码 $\delta\mu$ 稳定几何，第二阶段 $F_2$ 以更新后的均值为条件解码 $\delta s, \delta\alpha, \delta c, \delta r$：

\begin{align}
[\delta\mu] &= F_1(\mu^s, s^s, \alpha^s, c^s, r^s;\, f_\theta),\\
[\delta s, \delta\alpha, \delta c, \delta r] &= F_2(\mu^s + \delta\mu, s^s, \alpha^s, c^s, r^s;\, f_\theta).
\end{align}

为稳定训练初期，$F$ 末尾使用零初始化线性层（zero linear），确保 $P_\theta$ 初始输出为零，提供更有效的初始梯度。最终的编辑结果为 $X^r = X^s + \Delta$，$P_\theta$ 的训练损失为 $\mathcal{L}_{\text{distill}} = \mathbb{E}_{X^s, y, z}\mathbb{E}_{X^r}\left[d\left(R(X^r),\, x^e\right)\right]$，其中 $R$ 为可微渲染、$x^e$ 为 2D 编辑图像、$d$ 为 MSE 距离度量。VF-Editor 的统一设计使其能从多种 2D 编辑模型和策略蒸馏知识到单一预测器中，推理无需逐场景优化。

\paragraph{DisCo3D: 3D 到 2D 的逆向蒸馏}

与 VF-Editor 的"2D→3D"蒸馏方向相反，DisCo3D（2025.08）~\cite{2508_01684}采取了"3D→2D"的逆向蒸馏路线——从 3D NVS（Novel View Synthesis，新视角合成）扩散模型中提取跨视角一致性先验，注入 2D 编辑器以约束其多视角输出。

\textbf{动机。}VF-Editor 需大规模离线蒸馏且推理时完全绕过 2D 编辑器；DisCo3D 保留 2D 编辑器的灵活性，通过将 3D 一致性先验迁移到 2D 编辑器中来约束其多视角输出，但需要为每个场景微调 NVS 教师。具体而言，DisCo3D 先通过 LoRA 微调 3D U-Net NVS 模型的时间注意力层捕获源场景的跨视角一致性，训练目标为 $\mathcal{L}_{temporal} = \mathbb{E}_{t,\epsilon}\|\epsilon - \epsilon_\theta(z_t,t,c)\|^2$；再以微调后的 NVS 模型为教师，通过一致性蒸馏框架（基于 KL 散度最小化的蒸馏损失加正则化损失）将 3D 一致性先验迁移到 2D 编辑器的自注意力层 LoRA 模块中，保持其单视角编辑能力的同时约束多视角输出一致。这一设计的关键权衡在于：VF-Editor 训练通用预测器一次、推理即用，但需要大规模离线蒸馏；DisCo3D 保留 2D 编辑器的灵活性，但需要为每个场景微调 NVS 教师。

\paragraph{REALM: 开放世界推理分割定位}

上述方法聚焦于编辑操作本身的架构设计，而 REALM（2025.10）~\cite{2510_16410}转向编辑的前提环节——语义分割定位。

\textbf{动机。}精确编辑的前提是精确的 3D 区域定位，现有分割方法依赖 2D 掩码蒸馏或逐场景优化，缺乏开放世界的推理能力；REALM 利用 MLLM 推理能力进行开放世界分割，无需 3D 特定后训练。与 VF-Editor 直接训练变化预测器不同，REALM 利用 MLLM 推理能力在 3DGS 上进行开放世界推理分割，无需 3D 特定后训练，通过为每个高斯基元分配实例特征 $f_i \in \mathbb{R}^D$ 并训练分类器在渲染特征图上直接预测像素级实例 ID，为后续编辑提供精确的 3D 区域定位。

\paragraph{TransSplat: 最优传输语义编辑}

从 REALM 的分割定位出发，TransSplat（2026.04）~\cite{2604_19571}将语言驱动 3DGS 编辑重新建模为多视角不平衡语义传输问题，与 ProGDF 的连续变形场和 VF-Editor 的前馈预测不同，TransSplat 将编辑视为从 2D 编辑证据到 3D 高斯的显式最优传输分配。其核心流程为：首先通过置信度感知聚类将每个视角的密集编辑证据压缩为 $M_v$ 个提示感知原型 $p_m^v = (u_m^v, e_m^v, b_m^v)$，分别编码注意力加权质心、聚合语义特征和归一化目标质量；然后在每个视角内构建从可见高斯到原型的传输代价
\begin{equation}
C_{im}^v = \lambda_{geo}\|\pi_v(x_i) - u_m^v\|_2^2 + \lambda_{sem}(1 - \frac{s_i \cdot e_m^v}{\|s_i\|\|e_m^v\| + \epsilon}) + \lambda_{app}\|\phi_i^v - \phi_m^v\|_2^2,
\end{equation}

其中几何项约束图像平面局部性，语义项衡量高斯语义潜变量与原型语义特征的一致性，外观项通过局部视觉描述子抑制空间邻近但视觉内容不匹配的误对应。基于此代价求解熵正则化不平衡传输问题 $T_v^* = \arg\min_{T \geq 0} \langle T, C_v \rangle + \epsilon KL$，允许源-目标质量不匹配以适应编辑中可能引入的新语义或移除的旧结构。传输后每个高斯获得视角局部语义目标 $y_i^v$ 和传输置信度 $w_i^v$，再通过正则化重心问题恢复规范 3D 编辑场 $z_i^* = \arg\min_{z} \sum_{v \in V(i)} \lambda_i^v\|z - y_i^v\|_2^2 + \|z - s_i\|_2^2$，加权聚合多视角证据的同时向当前语义潜变量正则化以防止过度漂移。

\paragraph{FluSplat: 稀疏双视角前馈编辑}

沿 VF-Editor 的前馈推理路线进一步推进，FluSplat（2026.04）~\cite{2604_20038}将前馈范式从密集多视角扩展到稀疏双视角场景：仅需两个输入视角和文本编辑指令，即可通过单次前向传播直接生成编辑后的 3DGS 表示，完全消除逐场景优化。

\textbf{动机。}VF-Editor 依赖密集多视角输入，但实际场景中高质量多视角数据往往难以获取；FluSplat 将前馈范式压缩到双视角，通过跨视角一致性约束在图像域而非 3D 域保证编辑一致性。与 VF-Editor 在 3D 域训练变化预测器不同，FluSplat 将编辑一致性约束从 3D 域转移到图像域——基于 FLUX Kontext 整流流模型在秩 $r=4$ 的 LoRA 适配下微调跨视角一致性约束（包括全局扩散特征损失 GDFL 和局部编辑特征损失 LEFL），将编辑和重建解耦为跨视角一致图像编辑与基于 transformer 的前馈 3D 高斯重建两个阶段。

\paragraph{RoMaP: 部分级精细编辑}

前面方法多聚焦于全局或语义级编辑，而 RoMaP（2025.07，Robust Part-level Manipulation）~\cite{2507_11061}转向更精细的部分级编辑，需解决两大难题：如何在 3DGS 中准确定位目标部件，以及如何让 SDS 损失仅作用于该部件。

\textbf{动机。}物体级编辑的粒度太粗，无法实现"只改变杯子的把手"或"只修改椅子的靠背"等精细操作；RoMaP 通过 3D 几何感知标签预测和正则化 SDS 损失实现部件级精确编辑。在 3D 几何感知标签预测（3D-GALP, 3D Geometry-Aware Label Prediction）方面，为每个高斯引入球谐系数标签参数 $r_i$，将渲染标签图 $r^\theta = SH(r, \theta)$ 视为分割图，通过 $L_1(A(C^\theta), R^\theta)$ 损失与 2D 注意力图对齐。由于边界处高斯同时属于多个部分、不同视角标注不一致，RoMaP 定义标签软度 $S_i = H_i \cdot v_i$（其中 $H_i$ 为标签熵、$v_i$ 为 $r^\theta$ 跨视角方差），据此从高软度和低软度区域各选 $K/2$ 个锚点，通过近邻一致性损失 $L_{GALP} = \sum_{i \in \mathcal{A}}\sum_{j \in N_K(i)}\|r_j - r_i\|_1$ 精细化边界分割。在正则化 SDS 损失方面，RoMaP 提出组合损失
\begin{equation}
\mathcal{L}_{R\text{-}SDS} = \lambda_1 \mathcal{L}_{SDS}(c^\theta_{pr}, p_{edit}) + \lambda_2 \mathcal{M}(c^\theta_{pr}, \Phi(c^\theta_{pr})),
\end{equation}

其中第一项为高斯先验去除后的 SDS 损失（用中性颜色替换主导色彩先验 $c_{pr}$ 以避免对原始外观的固着），$\mathcal{M}$ 为利用 $M_{3D}$ 和 $M_{2D}$ 掩码限制修改仅作用于目标区域的掩码损失，$\Phi$ 为 SLaMP（Scheduled Latent Mixing and Part）编辑策略生成的锚定图像。SLaMP 通过时间相关混合系数 $F_t$ 在目标潜变量 $z_t$ 和原始潜变量 $z_{t,orig}$ 之间调度：$z_{t+1} = z_t(1 - F_t(1-M_{2D})) + z_{t,orig}F_t(1-M_{2D})$，前期低 $F_t$ 允许生成新内容，在切换步 $t_s$ 处急剧提高以保留原始图像上下文，平衡局部修改与全局身份保持。

\paragraph{TRACE: mesh 引导与视频扩散编辑}

与前面方法主要在 3DGS 基元层面操作不同，TRACE（2026.04）~\cite{2604_01207}引入 mesh 几何引导和视频扩散两条互补路线，实现高保真细粒度编辑。

\textbf{动机。}纯高斯操作缺乏结构化的几何约束，在大范围编辑中容易产生几何退化；TRACE 引入 mesh 作为几何代理提供结构约束，并通过视频扩散模型提供时序一致的编辑参考。其三阶段流水线首先通过多视角 3D 锚点合成建立视觉参考——利用水平拼接策略 $I_{ctx} = [I_1^* \ I_t]$ 以首帧编辑结果 $I_1^*$ 为视觉锚点，VLM 识别物理接触点生成空间提示 $y_{iter}$ 约束跨视角几何关系，并在新构建的 MV-TRACE 数据集上通过球面密集采样 $(r,\theta,\phi)$ 和大角度差异视角对 $(I_a, I_b)$ 微调 3D-LoRA。然后通过可触及几何对齐（TGA）将 LRM 生成的 mesh 与场景对齐——第一阶段求解 Procrustes 式相似变换 $\arg\min_{s,R,t}\sum_i\|P_{Mono}^i - (sRP_{Mesh}^i + t)\|^2$ 锁定全局旋转，第二阶段通过可微细化损失 $\mathcal{L}_{refine} = \lambda_{geo}\mathcal{L}_{geo} + \lambda_{mask}\mathcal{L}_{mask} + \mathbb{1}_{SDF}(\lambda_{sdf}\mathcal{L}_{SDF} + R_{reg})$ 从粗到精进行 7-DoF 参数优化。最后通过上下文视频掩码模块（CVM）利用球面线性插值（SLERP）自适应采样轨迹并执行上下文掩码修复，将编辑帧反向投影回统一的 3DGS 表示完成重建。

\paragraph{GaussEdit 与 SVGS: 编辑流程优化}

上述架构创新之外，还有若干方法从不同维度丰富文本引导编辑的能力。在\textbf{编辑流程优化}方向，GaussEdit（2025.09）~\cite{2509_26055}提出三阶段框架（区域选择→全局-局部优化→纹理增强），通过包围框 ROI 定义和自适应 SD/MVDream 交替策略实现精确的区域级编辑。SVGS（2026.03）~\cite{2603_28126}则反其道而行，仅从单视角编辑出发，利用多视角扩散模型将编辑传播到其他视角，降低了多视角协调的复杂度。

\paragraph{CP-GS 与 Invert3D: 个性化编辑}

在\textbf{个性化编辑}方向，CP-GS（2025.05）~\cite{2505_14537}通过预训练 image-to-3D 模型和迭代 LoRA 微调实现单图参考的 3DGS 个性化。Invert3D（2025.08）~\cite{2508_16932}首次建立 3DGS 与文本嵌入空间的对齐，通过相机条件化逆映射实现自然语言引导的 3D 内容个性化编辑~\cite{2505_11934,2507_04961,2508_17645,2512_20129}。这些方法共同展示了文本引导编辑从"修改现有场景"向"个性化定制场景"的演进趋势。

\subsubsection{交互式物体操控}

本路线介绍通过点击、框选、拖拽等直接交互方式操控 3DGS 场景中物体的方法。与文本引导编辑的"自动理解"模式不同，交互式操控追求"用户指哪里改哪里"的直接控制体验。3DGS 的显式表示为交互式编辑提供了天然优势：用户可以直接在渲染画面上点击、框选、拖拽，系统将 2D 操作映射到对应的 3D 高斯基元。

Point'n Move（2023.11）~\cite{2311_16737}是首个在 3DGS 上实现端到端交互式物体操作的框架。

\textbf{动机。}文本引导方法依赖语言理解来指定编辑目标，语义模糊时难以精确控制；3DGS 的显式点云表示为交互式编辑提供了天然优势——用户可直接在渲染画面上点击，系统将 2D 操作映射到对应的 3D 高斯基元。Point'n Move 的根本区别在于：用户只需在渲染画面上点击若干 2D 点，系统即可将其映射为精确的 3D 物体分割，进而实时执行平移、旋转、删除等操作。这直觉得益于 3DGS 的显式点云表示——每个高斯基元都是可直接操作的离散实体，无需通过隐式场反向传播来间接编辑。

\textbf{双阶段自提示分割。} Point'n Move 的核心是双阶段自提示分割算法（dual-stage self-prompting segmentation），充分利用了 3DGS 可将任意属性视为颜色进行泼溅渲染的特性。粗阶段为每个高斯球添加可训练的分割分数属性（segmentation score），将该分数视作颜色通过泼溅渲染生成 2D 分割图，再从中提取提示点送入 SAM（2023）~\cite{sam} 生成伪真值掩码，以分割损失
\begin{equation}
\mathcal{L}_{\text{coarse}} = L_m(M_{\text{sam}}, M_r)
\end{equation}

迭代优化（$M_{\text{sam}}$ 为 SAM 伪真值、$M_r$ 为渲染掩码），使用 SGD 优化器、学习率 1.0。这一跨视角自提示过程在所有提供的相机位姿上迭代传播，逐步将 2D 标注扩散为一致的 3D 分割。

精化阶段引入对偶分割分数（dual segmentation score）捕获不应被选中的区域，以弥补粗阶段的遗漏。首先以粗阶段高分高斯计算 3D 包围盒，盒外高斯的对偶分数初始化为接近 1、盒内初始化为 0；然后在全部视角上以联合损失继续训练：

\begin{equation}
\mathcal{L}_{\text{fine}} = \mathcal{L}_{\text{coarse}} + L_m(M_{sd}, M_{rd}) + \mathcal{L}_{\text{dual}},
\end{equation}

其中 $L_m(M_{sd}, M_{rd})$ 为对偶掩码损失（权重$\lambda_{dd}=0.1$），$\mathcal{L}_{\text{dual}} = L_m(-M_r, M_{rd}) + \lambda_{dd} L_m(-M_{rd}, M_r)$ 鼓励正掩码与对偶掩码不重叠。最终以拒绝-包含方案（rejection-inclusion）合并两层分数：除了高分割分数的高斯，包围盒内对偶分数低于阈值的高斯也被纳入，从而回收隐藏在物体表面下方、难以通过渲染梯度训练的高斯。

\textbf{操控与暴露修复。} 分割完成后，场景被分解为目标 $R_{obj}$ 与背景 $R_{bg}$，用户可实时执行平移、旋转、删除操作——只需变换高斯基元的位置 $\boldsymbol{\mu}$ 和旋转四元数 $r$，无需逐编辑训练。删除操作会在 $R_{bg}$ 中留下暴露区域。修复策略采用场景内容揭示剪枝（scene-content revealing pruning）：先移除远离未遮挡高斯的选中高斯，最小化修复范围；再用 2D 修复器（LaMa）分别修复 RGB 图像和深度图；最后通过重投影将修复区域初始化为新高斯，以掩码外像素 $L_1$ + SSIM 颜色损失 $\mathcal{L}_{\text{ocolor}}$（$\lambda_{\text{ssim}}=0.2$）、深度 $L_1$ 损失 $\mathcal{L}_{\text{depth}}$（$\lambda_{\text{depth}}=1.0$）和掩码内 LPIPS 感知损失 $\mathcal{L}_{\text{icolor}}$（$\lambda_{\text{lpips}}=1.0$）微调 $R_{bg}$ 完成修复。感知损失的使用是因为 2D 修复结果不保证多视角一致，严格像素级损失会损害质量。

Point'n Move 的点击分割虽然直观，但本质上是一种几何驱动的低层分割——用户必须在画面上精确点击目标物体，无法通过语言描述（"选中那个红色的杯子"）来指定编辑对象。后续工作从两个不同方向突破了这一限制：一个走向语言驱动的语义解耦，另一个走向结构化基元的部件级解耦。

\textbf{语言驱动解耦：3DitScene。} 3DitScene（2024.05）~\cite{2405_18424}保留了 Point'n Move 的对象级操控能力，但将分割方式从"用户点击 → SAM 分割"替换为"语言/框选 → CLIP 语义匹配"，使用户无需精确点击即可选中目标。具体地，3DitScene 为每个 3D 高斯基元额外嵌入一个 CLIP 语言向量 $e \in \mathbb{R}^C$，利用 3DGS 的可微分泼溅将该向量渲染为 2D 语义特征图 $E \in \mathbb{R}^{C \times H \times W}$，以蒸馏损失
\begin{equation}
\mathcal{L}_{\text{distill}} = \sum_i \|(E - g(I \odot M_i)) \odot M_i\|_2^2
\end{equation}

对齐 SAM 分割掩码 $M_i$ 区域内的 CLIP 图像编码 $g(\cdot)$。为降低高维特征向量的显存开销，3DitScene 额外训练一个自编码器压缩嵌入空间（遵循 LangSplat（2023.12）~\cite{Qin2023LangSplat} 的策略）。蒸馏完成后，用户可通过两种方式查询对象：（1）文本提示——计算语言嵌入 $e$ 与文本提示嵌入 $e_l$ 的相关性分数 $\text{score} = \frac{e \cdot e_l}{e_l \cdot e_l + e \cdot e_{\text{canon}}}$（其中 $e_{\text{canon}}$ 为"object/things/stuff/texture"等基准短语的 CLIP 嵌入，用于抑制背景噪声），低于阈值的高斯被排除，余下即为目标对象；（2）包围框——先粗选框内高斯，再以其语言嵌入做 K-Means 聚类，丢弃占比较小的簇，保留主簇为选中对象。解耦后的高斯支持平移、旋转、删除、风格化等单对象编辑，也可在训练阶段对场景布局做随机增广（随机平移、旋转、移除对象），显著改善遮挡区域的渲染质量。

\textbf{部件级解耦：GaussianBlock。} 如果说 3DitScene 用语义特征替代了 Point'n Move 的点击提示，GaussianBlock（2024.10）~\cite{2410_01535}则更进一步——用参数化超二次曲面（Superquadric）替代了 Point'n Move 的无结构高斯集合，从根本上赋予场景以部件级的结构化表示。

\textbf{动机。}3DitScene 的语义特征虽然实现了语言驱动选择，但场景底层仍是无结构高斯集合，无法像积木一样拖拽、旋转部件；GaussianBlock 用参数化超二次曲面提供结构化表示，使用户可像操作积木一样操控语义部件。GaussianBlock 采用两阶段混合表示策略：第一阶段以参数化超二次曲面为基元，每个超二次曲面仅用少量连续参数（旋转 $r \in \mathbb{R}^6$、平移 $t \in \mathbb{R}^3$、缩放 $s \in \mathbb{R}^3$、形状 $\epsilon \in \mathbb{R}^2$）即可表达丰富的几何形状；通过软双光栅化（soft dual-rasterization）将超二次曲面顶点投影为 2D 点提示和包围框提示，查询预训练 SAM 获取注意力图 $A_k$；随后用 HDBSCAN 对 $A_k$ 聚类，以注意力引导中心化损失（AC loss）将离群顶点推向主聚类质心，实现语义一致的超二次曲面解耦；动态分裂-融合策略进一步防止局部最优（分裂阈值设为包围框对角线的 $0.4$ 倍，新超二次曲面缩放继承为 $s_z = 0.4s_k$）。第二阶段将 3D 高斯绑定到超二次曲面三角面片上，通过局部参数化在光栅化时映射到世界空间，配合位置正则化 $L_{pos} = \|\max(\|\cdot\|_2, \tau_{pos}) - \tau_{pos}\|_2$（$\tau_{pos}=0.5$）保持二者耦合。最终用户可像操作积木一样拖拽、旋转、缩放语义部件。

近期交互式方法在实用性和灵活性上持续进步。iSegMan（2025.05）~\cite{2505_11934}提出首个无需场景特定训练的交互式 3DGS 分割与操控框架，利用极线约束传播（EIP）和可见性投票（VGV）实现高效精确的区域选择。SplatPainter（2025.12）~\cite{2512_05354}通过紧凑特征丰富的 3DGS 表示和 Test-Time Training 层实现交互式迭代精修，支持局部精修、绘画和全局重着色。

在颜色编辑方面，ReCoGS（2025.11）~\cite{2511_18441}实现预训练 3DGS 场景的实时交互式重着色，通过 2D 选择掩码反投影为 3D 点云再投射到训练视角图像实现精确区域选择。

\textbf{动机。}前述方法多聚焦于形状或纹理编辑，颜色编辑缺乏实时交互能力；ReCoGS 通过 2D 掩码反投影到 3D 再投射到训练视角，实现精确区域选择和实时重着色。VIRGi（2026.03）~\cite{2603_02986}是首个面向 3DGS 的光感逼真外观编辑方法，其关键设计是漫反射/镜面反射双 MLP 解耦架构：漫反射 MLP $M_{df}$ 仅接收 hashgrid 特征 $x$（不含视角信息，因此无法建模视角相关效果），镜面反射 MLP $M_{sp}$ 同时接收 $x$ 和视角方向 $d$，并通过单向残差连接（$M_{df}$ 第一层→$M_{sp}$ 第二层，$M_{df}$ 第二层→$M_{sp}$ 第三层）让镜面反射分支获得漫反射的上下文信息，最终颜色由 $\sigma(C_{df} + C_{sp})$ 计算。训练损失 $L = \|C_n - RGB_{GT}\|_{L_1} + \lambda_{sp} \|C_s\|^2_2$ 中的镜面反射惩罚系数 $\lambda_{sp}$ 从 $0.25$ 线性衰减至 $0.05$，迫使网络优先用漫反射分支解释颜色变化。在多视角训练策略中，每个批次从 5 张不同视角的图像采样瓦片（tile），使同一高斯同时接收多视角梯度，隐式促进漫反射/镜面反射的有效分离。编辑阶段仅需用户修改一张视角图像，系统微调漫反射 MLP 最后一层和分割 MLP，即可在 2 秒内将颜色编辑传播到整个场景。ColorGradedGaussians（2026.04）~\cite{2604_01551}实现专业级实时调色工作流——调色板编辑、色调曲线、像素级约束三者协同。

交互式操控还有多种拓展方向。在\textbf{参考图引导}方面，ICE-G（2024.06）~\cite{2406_08488}通过 SAM 分割和 DINO（自监督视觉 Transformer）语义匹配实现单张参考图像编辑。MiraGe（2024.10）~\cite{2410_01521}利用镜面反射将 2D 图像映射到 3D 空间。在\textbf{特定创作形式}方面，SandSim（2026.04）~\cite{2604_27572}通过曲线引导高斯表示从单张图像重建沙画创作过程。其核心做法是将 3DGS 的线性结构用曲线参数化，从单张图像中恢复沙画艺术家的创作轨迹。在\textbf{效率与专用工具}方面，Instant Colorization（2026.04）~\cite{2604_17155}用正规方程实现高效 2D 到 3D 信息映射（比梯度下降快一个数量级）。SketchFaceGS（2026.04）~\cite{2604_19202}是首个基于草图的实时 3D 高斯头部生成与编辑框架~\cite{2510_09537,2511_22228,2601_03319,2603_28152,2603_15368}。

\begin{takeaway}
语义与文本引导编辑通过语言理解桥接用户意图与3D场景操控。GaussianBlock 用参数化超二次曲面实现部件级"积木式"操控，VIRGi 的漫反射/镜面反射双 MLP 解耦架构在 2 秒内将颜色编辑传播到整个场景。核心趋势是从交互式点击选择走向语言驱动的语义级编辑，从形状/纹理编辑扩展到物理感知的外观编辑。
\end{takeaway}

%=====================================================================


\section{风格化与场景修复}
\label{sec:09_2}

本节讨论两个紧密相关的编辑方向：风格化与艺术编辑（第~\ref{sec:09_2}~节），将艺术风格迁移到 3DGS 场景中；场景修复与物体移除，在多视角一致的前提下恢复被遮挡的背景内容。

\subsection{外观风格化与场景修复}
\label{ss:09_2}

本节讨论两条技术路线：（1）风格化与艺术编辑；（2）场景修复与物体移除。

\subsubsection{风格化与艺术编辑}

本路线介绍将艺术风格迁移到 3DGS 场景中并保持多视角一致性的方法。风格化将一幅参考图像的艺术特征（如梵高的笔触、毕加索的色块）迁移到 3DGS 场景中，同时保持场景的多视角一致性和几何完整性。这一任务的核心矛盾在于：离散高斯的外观-几何耦合使纹理迁移变得困难——改变高斯颜色会同时影响它在所有视角下的外观，而风格化要求不同表面区域呈现不同的纹理细节。

\begin{figure}[htbp]
\centering
\resizebox{0.95\textwidth}{!}{\begin{tikzpicture}[
  event/.style={rectangle, rounded corners=2pt, draw=#1, fill=#1!10, line width=0.5pt, minimum height=0.6cm, align=center, font=\scriptsize, inner sep=2pt},
  milestone/.style={rectangle, rounded corners=3pt, draw=#1, fill=#1!20, line width=0.8pt, minimum height=0.65cm, align=center, font=\scriptsize\bfseries, inner sep=3pt},
]
\draw[thick, -{stealth[length=5pt]}, draw=techgray] (0, 0) -- (14, 0)
  node[right, font=\scriptsize] {时间};

% Year markers
\foreach \x/\y in {1/2024.0, 4/2024.5, 7/2025.0, 10/2025.5, 13/2026.0} {
  \draw[techgray] (\x, -0.1) -- (\x, 0.1);
  \node[below, font=\scriptsize, text=techgray] at (\x, -0.15) {\y};
}

% Phase bands
\fill[techblue!5] (0.5, 0.3) rectangle (3.5, 3.2);
\fill[techgreen!5] (3.5, 0.3) rectangle (7.5, 3.2);
\fill[techorange!5] (7.5, 0.3) rectangle (10.5, 3.2);
\fill[techpurple!5] (10.5, 0.3) rectangle (13.5, 3.2);

\node[font=\scriptsize\bfseries, text=techblue] at (2, 3.4) {优化式};
\node[font=\scriptsize\bfseries, text=techgreen] at (5.5, 3.4) {分布匹配};
\node[font=\scriptsize\bfseries, text=techorange] at (9, 3.4) {零样本/前馈};
\node[font=\scriptsize\bfseries, text=techpurple] at (12, 3.4) {统一框架};

% Events - Optimization era (3 rows: y=0.7, 1.5, 2.3)
\node[milestone=techblue] at (1.5, 2.3) {StylizedGS};
\node[event=techblue] at (2.8, 1.5) {StyleSplat, G-Style};
\node[event=techblue] at (1.5, 0.7) {InstantStyleGS};

% Distribution matching
\node[milestone=techgreen] at (4.5, 2.3) {WaSt-3D};
\node[event=techgreen] at (5.8, 1.5) {ReGS, ArtNVG};
\node[event=techgreen] at (4.8, 0.7) {SGSST};

% Zero-shot / Feed-forward
\node[milestone=techorange] at (8.5, 2.3) {FPGS, 4DStyleGS};
\node[event=techorange] at (9.5, 1.5) {Stylos};
\node[event=techorange] at (8.5, 0.7) {AnyStyle};

% Unified frameworks
\node[milestone=techpurple] at (12, 2.3) {CLIPGaussian};
\node[event=techpurple] at (12, 1.5) {SSGaussian, M2StyleGS};
\node[event=techpurple] at (12, 0.7) {ABC-GS};

\end{tikzpicture}
}
\caption{3DGS 风格化方法的四个发展阶段。从优化式（逐场景迭代，质量高但速度慢）到分布匹配（理论突破），再到零样本/前馈（一次训练任意风格），最终走向统一多模态框架（CLIPGaussian 同时支持 2D/视频/3D/4D）。}
\label{fig:ch09_stylization_timeline}
\end{figure}

\paragraph{StylizedGS: 场景级 NNFM 风格化}

最早的 3DGS 风格化方法沿袭神经风格迁移的优化范式，通过迭代损失函数将风格特征注入高斯参数。

\textbf{动机。}传统 Gram 矩阵风格迁移倾向于丢失局部纹理结构，因为 Gram 矩阵仅捕获全局二阶统计量；NNFM（最近邻特征匹配）通过逐像素搜索风格图中最相似的特征点，能更精确地保留风格的局部纹理。

StylizedGS（2024.04）~\cite{2404_05220}是首个基于 3DGS 的可控场景风格化框架。传统神经风格迁移（如 Gatys 等）依赖 Gram 矩阵捕获特征图的全局二阶统计量进行风格匹配，这种方式倾向于丢失局部纹理结构。StylizedGS 的核心改动是用最近邻特征匹配损失（NNFM）替代 Gram 矩阵损失——直觉上，即为渲染图像的每个像素在风格图中找到"最像"的特征点，然后拉近二者距离：

\begin{equation}
\mathcal{L}_{\text{style}} = \frac{1}{N} \sum_{i,j} \min_{i',j'} D\!\bigl(F_{\text{render}}(i,j),\; F_{\text{style}}(i',j')\bigr)
\end{equation}

其中 $D(\cdot,\cdot)$ 为余弦距离，$(i^*,j^*) = \arg\min_{i',j'} D(F_{\text{render}}(i,j), F_{\text{style}}(i',j'))$ 为最近邻位置。相比 Gram 矩阵匹配全局统计量，NNFM 逐像素搜索最近邻的方式能更精确地保留风格图像的局部纹理结构。为防止过度风格化破坏场景内容，还引入内容损失 $\mathcal{L}_{\text{content}} = \frac{1}{HW}\|F_{\text{content}} - F_{\text{render}}\|^2$，约束渲染特征不偏离原始内容特征过远。

在几何保持方面，StylizedGS 引入深度保持损失 $\mathcal{L}_{\text{depth}} = \frac{1}{HW}\|D_{\text{origin}} - D_{\text{render}}\|^2$，利用 3DGS 自身的 $\alpha$-混合 深度图 $\hat{d} = \sum_i T_i \alpha_i d_i$ 作为锚定参考，无需额外深度估计器。此外还加入尺度正则 $\|\Delta \mathbf{s}\|$、不透明度正则 $\|\Delta \alpha\|$ 和 2D 域全变分正则 $\mathcal{L}_{\text{tv}}$，风格化阶段总损失为六项加权和：$\mathcal{L} = \lambda_{\text{sty}}\mathcal{L}_{\text{style}} + \lambda_{\text{con}}\mathcal{L}_{\text{content}} + \lambda_{\text{dep}}\mathcal{L}_{\text{depth}} + \lambda_{\text{sca}}\mathcal{L}_{\text{reg}}^{\Delta s} + \lambda_{\text{opa}}\mathcal{L}_{\text{reg}}^{\Delta \alpha} + \lambda_{\text{tv}}\mathcal{L}_{\text{tv}}$。

在感知控制方面，该框架支持三个维度的独立调控：颜色控制通过在 YIQ 色彩空间（将亮度 Y 与色度 I/Q 解耦）中仅对亮度通道计算 NNFM 损失，允许用户保留原始色调或混搭不同风格的配色；尺度控制通过选择 VGG 不同层的特征图及调整各层权重来调节笔触粗细；空间区域控制允许用户通过点交互（SAM）或语言提示（LangSAM）为场景的不同区域指定不同风格。

\paragraph{StyleSplat: 物体级风格化}

与 StylizedGS 的场景级全局风格化不同，StyleSplat（2024.07）~\cite{2407_09473}将风格化粒度推进到物体级——场景中的不同物体可以被赋予不同的艺术风格，而这正是场景级方法无法实现的核心局限。

\textbf{动机。}场景级方法对所有区域施加统一风格，无法让场景中不同物体呈现不同艺术效果；StyleSplat 通过 Gaussian Grouping 机制先分割再独立风格化，实现场景内多物体独立风格化。其关键设计是 Gaussian Grouping 机制：为每个 3D 高斯分配一个 16 维可学习特征向量 $\mathbf{e}_i \in \mathbb{R}^{16}$，该编码与球谐系数同步优化，渲染时通过 $\alpha$-混合聚合为像素级特征：

\begin{equation}
E_{\text{id}} = \sum_{i}^{N} e_i' \alpha_i' \prod_{j=1}^{i-1}(1-\alpha_j')
\end{equation}

渲染特征经分类器输出语义标签，以交叉熵损失与 DEVA+SAM 生成的时序一致 2D 分割掩码对齐，辅以空间一致性损失确保近邻高斯特征相似。分割完成后，仅解冻用户指定物体的球谐系数进行梯度更新，其余高斯属性全部冻结，风格化被严格约束在目标物体内部。NNFM 损失仅对选物体的 VGG 特征计算，实现首个支持场景内多物体独立风格化的轻量方法。

\paragraph{G-Style: 多尺度损失与自适应几何}

上述方法（StylizedGS、StyleSplat）在风格化过程中都严格保持高斯几何不变——StylizedGS 的深度保持损失、尺度正则和不透明度正则本质上将位置与形状锁定在原始重建值，只修改颜色（球谐系数）；StyleSplat 进一步冻结非目标物体的全部属性。

\textbf{动机。}锁定几何意味着欠采样区域的粗糙分辨率被继承；此外 NNFM 仅捕获高频纹理，无法建模风格图像中的大尺度构图与色彩布局。G-Style 从这两个局限出发，允许风格化过程中高斯几何动态调整，并引入多尺度风格损失。

G-Style（2024.08）~\cite{2408_15695}正是从这两个局限出发，做出两项关键扩展。

\textbf{第一项扩展——多尺度风格损失}：在 StylizedGS 的 NNFM 高频损失基础上，新增 CLIP 图像编码器计算低频风格距离 $\mathcal{L}_{\text{CLIP}} = \frac{1}{K}\sum_k \|C(I_r^k) - C(I_s)\|_2^2$，以捕获 NNFM 难以覆盖的大尺度图案和整体构图。高频风格仍沿用 VGG-16 第三块特征图的 NNFM 损失捕获笔触等精细纹理，辅以内容损失 $\mathcal{L}_C$ 和全变分正则 $\mathcal{L}_{TV}$，总损失为 $\mathcal{L} = \lambda_{\text{CLIP}}\mathcal{L}_{\text{CLIP}} + \lambda_{\text{NNFM}}\mathcal{L}_{\text{NNFM}} + \lambda_C\mathcal{L}_C + \lambda_{TV}\mathcal{L}_{TV}$（前向场景 $\lambda_{\text{NNFM}}{=}100$，360° 场景降为 10，以适配后者的多视角一致性需求）。此外，在风格化开始前，G-Style 还执行一步颜色匹配预处理：解析求解线性变换 $A$，使内容场景像素颜色的均值和协方差矩阵与风格图像对齐（$\mathbb{E}[AC]=\mathbb{E}[S]$，$\text{Cov}(AC)=\text{Cov}(S)$），为后续优化提供一个色彩分布一致的初始化。

\textbf{第二项扩展——自适应几何优化}：这是 G-Style 区别于此前所有 3DGS 风格化方法的核心——风格化过程中高斯的位置和尺度不再固定，而是根据风格需求动态调整。具体分为两个阶段。预处理阶段将过扁高斯（投影面积超过 $\mu_A + \sigma \cdot t_f$，$t_f$ 初始 1.1，每轮乘以 1.125）和过窄高斯（拉长因子 $E_i > 1.5$）进行归一化，使高斯尺寸分布均匀，为后续细节合成提供均匀的几何基础。风格化迭代中维护累积缓冲区 $B$ 存储风格颜色梯度 $\|\nabla c_s\|$ 的范数，周期性分裂梯度最高的高斯以在需要细节的区域自动增加高斯密度，分裂后重新优化位置 $\boldsymbol{\mu}$、尺度 $\mathbf{s}$、不透明度 $\alpha$ 和原始颜色 $c_{\text{gt}}$ 以恢复几何精度。为避免球谐系数引入视角依赖的风格伪影，G-Style 仅保留球谐的零阶项（即漫反射颜色），将颜色存储从 192 字节压缩至 12 字节，腾出空间容纳分裂产生的新高斯。整个风格化流程仅需 15 个 epoch 即可收敛。

\paragraph{InstantStyleGaussian: 简化风格化流程}

沿同一优化范式，InstantStyleGaussian（2024.08）~\cite{2408_04249}则从另一个维度简化了风格化流程：结合 InstantStyle 扩散模型和改进的迭代数据集更新策略，仅需目标风格图像即可快速完成风格编辑，无需手动调参损失权重。

\textbf{动机。}StylizedGS 和 StyleSplat 需要手动平衡风格损失、内容损失和正则化损失的权重，不同场景需要不同设置；InstantStyleGaussian 利用扩散模型自动处理风格-内容权衡，消除了手动调参。

\paragraph{分布匹配风格化}

前述优化式方法（如 StylizedGS、G-Style）本质上都在做逐像素或逐高斯的特征匹配——为每个内容点在风格图中找最近邻然后拉近。这种策略能捕获局部纹理，却忽略了全局分布结构：风格的"感觉"不只来自单个像素的颜色，更来自整体色彩和几何的统计分布。

\textbf{动机。}逐像素特征匹配丢失了风格的全局分布结构；WaSt-3D 将风格化重新定义为两个高斯集合之间的最优传输分布对齐，从"逐像素模仿"提升到"分布级别对齐"。WaSt-3D（2024.09，Wasserstein-3D）~\cite{2409_17917}提出了一个根本不同的视角：不做逐像素模仿，而是将风格化视为两个高斯集合之间的分布对齐问题——直觉上类似于拼贴艺术（assemblage）：用风格场景的各种碎片"组装"出内容场景的形状，正如 Arcimboldo 用水果蔬菜拼成人像。

这一分布对齐的核心工具是最优传输（Optimal Transport）。给定内容场景的高斯分布 $\mu$ 和风格场景的高斯分布 $\nu$，最优传输寻找一个"搬运计划" $\pi$，以最小代价将 $\mu$ 变形为 $\nu$。具体使用熵正则化 Wasserstein-2 距离：

\begin{equation}
  W_\varepsilon(\mu, \nu) = \min_{\pi \in \Pi(\mu,\nu)} \int c(x,y) \, d\pi(x,y) + \varepsilon \, \text{KL}(\pi \| \mu \otimes \nu),
\end{equation}

其中 $\Pi(\mu,\nu)$ 为所有满足边缘约束的联合分布（传输计划），$c(x,y)$ 为将高斯 $x$ 搬运到 $y$ 的代价（颜色与位置联合的欧氏距离），$\varepsilon > 0$ 为熵正则化系数。$\varepsilon$ 大时传输计划更平滑（每个风格高斯均匀地搬运到所有内容高斯），$\varepsilon$ 小时趋向一一对应（每个风格高斯只搬运到最近的内容高斯）。然而，$W_\varepsilon$ 本身不是度量（$W_\varepsilon(\mu,\mu) \neq 0$），因此采用其去偏版本 Sinkhorn 散度：

\begin{equation}
  \text{SD}_{2,\varepsilon}(\mu, \nu) = W_\varepsilon(\mu, \nu) - \tfrac{1}{2}W_\varepsilon(\mu, \mu) - \tfrac{1}{2}W_\varepsilon(\nu, \nu),
\end{equation}

当 $\varepsilon \to \infty$ 时退化为 MMD（最大均值差异），$\varepsilon \to 0$ 时退化为精确 $W_2$ 距离。Sinkhorn 散度满足度量性质且可通过 Sinkhorn 迭代高效计算，是实际优化的目标函数。

直接对整个场景计算最优传输不可行（高斯数量达百万级），因此 WaSt-3D 采用分治策略，流程分三步。

\textbf{第一步}，对内容场景做 K-Means 聚类，将高斯按颜色和位置分为不相交子集 $\{C_i\}$。聚类前先施加各向同性正则化 $\mathcal{L}_{\text{aniso}} = \frac{1}{|G|}\sum_{g \in G}\bigl(\frac{\max(g_S)}{\min(g_S)} - r\bigr)$ 强制高斯趋近球形，防止针状伪影在分割时产生碎片。

\textbf{第二步}，为每个内容簇 $C_i$ 寻找最佳匹配的风格片段。通过约束优化——仅允许平移 $t_i \in \mathbb{R}^3$、旋转 $R_i \in \text{SO}(3)$ 和缩放 $S_i \in \text{diag}(\mathbb{R}_+^3)$——将 $C_i$ 的坐标映射到风格场景坐标空间，取 $k$ 最近邻高斯构成匹配簇 $D_i(C_i)$。

\textbf{第三步}，对每对 $(C_i, D_i(C_i))$ 分别计算 Sinkhorn 散度并求和，作为最终优化目标：

\begin{equation}
  \mathcal{L}_{\text{opt}} = \sum_{i} \text{SD}_{2,\varepsilon}^2\bigl(C_i, D_i(C_i)\bigr).
\end{equation}

优化过程中同时调整风格簇高斯的颜色和位置坐标，实现几何与纹理的联合迁移。整个流程无需训练任何网络——这一视角转变将风格化从"逐像素外观模仿"提升到"分布级别的对齐"。

\paragraph{4DStyleGaussian 与 ZDySS: 零样本风格化先驱}

优化式方法需要针对每对内容-风格进行迭代优化，限制了实际应用的便利性。零样本（Zero-shot）和前馈（Feed-forward）方法追求"一次训练，任意风格"。

4DStyleGaussian（2024.10）~\cite{2410_10412}是首个零样本 4D 风格迁移框架，通过可逆神经网络和 4D 风格变换矩阵实现实时任意风格迁移，同时保持内容保真度、多视角一致性和时序连贯性。ZDySS（2025.01）~\cite{2501_03875}将风格迁移应用于学习到的特征向量而非渲染特征图，增强时空一致性。

\paragraph{FPGS: 前馈式逼真风格迁移}

直觉上，FPGS 类似于 Photoshop 的"匹配颜色"功能——只需一张风格参考图，就能将整个 3D 场景的色彩和纹理统计量一次性"刷"上去，无需逐场景迭代优化。

\textbf{动机。}优化式方法需要针对每对内容-风格进行分钟到小时级的迭代优化，严重限制实用性；FPGS 利用 3DGS 的显式表示特性，在渲染前一次性对所有高斯施加 AdaIN 风格化，将风格化与渲染解耦实现实时风格迁移。FPGS（2025.03）~\cite{2503_09635}是首个支持大规模 3DGS 场景前馈式逼真风格迁移的方法。此前的优化式方法需要针对每对内容-风格进行迭代优化，或对每个场景进行耗时的额外训练阶段（超过 10 小时），且仅能处理小规模场景。FPGS 的核心思想是将风格化能力直接嵌入 3D 高斯表示中，构建可风格化特征场（stylizable feature field）：场景内容场（scene content field）嵌入几何与外观内容特征，场景语义场（scene semantic field）嵌入 DINO 语义特征用于跨场景-参考图匹配。风格迁移通过自适应实例归一化（AdaIN）在前馈阶段一次性完成——对每个 3D 高斯的内容特征施加风格统计量：

\begin{equation}
  \text{AdaIN}(\mathbf{F}(c), \mathbf{F}(s)) = \boldsymbol{\sigma}_s \frac{\mathbf{F}(c) - \boldsymbol{\mu}_c}{\boldsymbol{\sigma}_c} + \boldsymbol{\mu}_s,
\end{equation}

其中 $\mathbf{F}(c)$ 和 $\mathbf{F}(s)$ 分别为内容和风格图像的 VGG 特征，$(\boldsymbol{\mu}_s, \boldsymbol{\sigma}_s)$ 为风格图像特征的通道均值与标准差，$(\boldsymbol{\mu}_c, \boldsymbol{\sigma}_c)$ 为内容特征的对应统计量。关键在于，利用 3DGS 的显式表示特性，FPGS 在渲染前一次性对整个场景的所有高斯施加 AdaIN，将风格化与渲染解耦，从而实现风格化后的实时渲染（而前序 NeRF 版本 FPRF 因隐式表示需要每次渲染时重新风格化）。对于多参考图场景，FPGS 将多张风格参考图的语义特征（DINO）和风格统计量（VGG 均值/标准差）聚类构建风格字典 $\mathcal{D}$，通过场景语义场与字典键的交叉相关矩阵计算语义加权风格码，实现语义感知的局部 AdaIN——场景中的天空区域自动匹配参考图中的天空风格，建筑匹配建筑风格。整个训练过程仅需约 24 分钟（30000 次迭代），风格化阶段无需任何额外优化。

\paragraph{Stylos: 双路径 Transformer 风格迁移}

沿前馈风格化路线，后续方法从架构和模态两个维度持续拓展。Stylos（2025.09）~\cite{2509_26455}设计双路径 Transformer 主干——几何预测路径保留自注意力保持几何保真度，风格路径通过全局交叉注意力注入确保视觉一致性，支持从单视角到多视角的零样本 3D 风格迁移，无需逐场景优化。

\paragraph{AnyStyle 与 M2StyleGS: 多模态前馈风格化}

沿前馈路线进一步推进，AnyStyle（2026.02）~\cite{2602_04043}基于 ControlNet 启发的零初始化卷积实现架构无关的风格注入，仅需轻量微调即可集成到预训练 3DGS 骨干网络中，支持文本和图像双模态控制。M2StyleGS（2026.04）~\cite{2604_03773}首次实现多模态（文本/图像）参考的实时 3DGS 风格迁移，提出 subdivisive flow 精确对齐 CLIP 与 VGG 特征空间。

\paragraph{Multi-StyleGS: 多风格场景风格化}

实际场景中，用户可能希望场景不同区域呈现不同风格，或同时利用文本和图像作为风格参考。

Multi-StyleGS（2025.06）~\cite{2506_06846}首次实现基于 3DGS 的多风格 3D 场景风格化，通过二部图匹配自动建立风格图像与场景区域的对应关系。

\paragraph{CLIPGaussian: 统一多模态风格框架}

在此之前，风格迁移方法高度碎片化：2D 图像、视频、3D 物体和 4D 场景各有专属方案，文本引导和图像引导需要不同模型，用户必须针对特定任务选择特定工具。

\textbf{动机。}风格迁移方法高度碎片化，不同模态和任务需要不同工具；CLIPGaussian 的核心洞察是：所有模态都可以用高斯基元统一表示，一旦统一，风格迁移就归结为同一操作——微调高斯参数。直觉上，CLIPGaussian 就像是给不同语言的翻译器装上同一个引擎——一旦所有内容都统一为高斯基元表示，风格迁移就变成同一套操作流程。CLIPGaussian（2025.05）~\cite{2505_22854}的关键洞察是：尽管这些模态的数据形式不同，它们都可以用高斯基元统一表示——2D 图像用 MiRaGe 的 2D 高斯、视频用 VeGaS 的时序高斯、3D 场景用经典 3DGS、4D 场景用 D-MiSo 的动态高斯。一旦统一为高斯表示，风格迁移就可以归结为同一操作：微调高斯参数。

具体而言，CLIPGaussian 采用两阶段流程。第一阶段用现有高斯模型训练场景，得到基座模型 $\mathcal{G}$。第二阶段冻结高斯数量，仅微调颜色（SH 系数）和几何参数，通过四项损失的加权组合驱动风格化：

\begin{equation}
  \mathcal{L}_{\text{total}} = \lambda_d \mathcal{L}_d + \lambda_p \mathcal{L}_p + \lambda_c \mathcal{L}_c + \lambda_b \mathcal{L}_b,
\end{equation}

其中方向 CLIP 损失 $\mathcal{L}_d$ 在全局尺度上将风格因子 $S$（文本提示或风格图像）的 CLIP 方向与重建图像的 CLIP 方向对齐；补丁 CLIP 损失 $\mathcal{L}_p$ 对随机裁剪的局部块施加同样的方向约束，确保风格细节渗透到局部纹理而非仅改变全局色调；VGG 内容损失 $\mathcal{L}_c$ 保持风格化结果与原始内容在 conv4\_2/conv5\_2 特征上的相似度；背景损失 $\mathcal{L}_b$ 通过 L1 距离约束背景像素不引入风格伪影。由于不改变高斯数量，风格化后的模型与原始模型体积相同，且支持通过线性插值高斯参数实现风格之间的平滑过渡。整个方法作为插件模块集成到任意高斯流程中，无需大型生成模型或从头重新训练。

\paragraph{ABC-GS: 分割掩码风格对齐}

沿多风格路线，后续方法从不同角度增强风格化能力。ABC-GS（2025.03）~\cite{2503_22218}通过分割掩码实现场景内容与风格特征的精确对齐。其 FAST Loss（Feature Alignment Style Transfer Loss）的核心思想是：对内容掩码 $M_c$ 和风格掩码 $M_s$ 中匹配的区域，计算 VGG 特征空间的 Gram 矩阵差异作为风格损失，同时通过高斯语义权重 $w_i^j = o_i(p) \cdot T_i^j(p) \cdot M_c^j(p)$ 将 2D 掩码反投影到 3D 高斯空间，确保风格只作用于目标区域。支持单图像、组合式和语义感知三种风格迁移类型。

\paragraph{SSGaussian 与 ReGS: 一致性与耦合问题}

沿多风格路线进一步推进。SSGaussian（2025.09）~\cite{2509_04379}首次在扩散模型中插入跨视角注意力确保多视角风格一致性，并实现实例级风格迁移。ReGS（2024.07）~\cite{2407_07220}针对参考图像风格化中离散高斯的外观-几何耦合问题，通过纹理引导控制机制自适应调整局部高斯的几何排布以捕捉连续纹理细节。

\paragraph{SGSST: 超高分辨率风格迁移}

除上述主线演进外，风格化领域还有多个值得关注的新方向。

在\textbf{分辨率与质量}方面，SGSST（2024.12）~\cite{2412_03371}提出 SOS 多尺度损失，首次实现超高分辨率（4K 以上）3DGS 场景风格迁移。

\paragraph{ArtNVG 与 StyleMe3D: 分离控制与高保真}

沿分辨率与质量方向继续推进。ArtNVG（2024.12）~\cite{2412_18783}通过内容-风格分离控制和邻域视角注意力对齐确保局部一致性；StyleMe3D（2025.04）~\cite{2504_15281}通过动态风格分数蒸馏、CLIP 双流对齐和 VGG 的 SOS 模块实现多层次高保真风格化。

\paragraph{Morpheus: 外观-形状解耦风格化}

在\textbf{外观-形状解耦}方面，Morpheus（2025.03）~\cite{2503_02009}通过 RGBD 扩散模型独立控制外观和形状的风格化强度，为用户提供更灵活的风格调节维度。

\paragraph{推理效率提升方法}

在推理效率方面，免优化方法（2025.08）~\cite{2508_05813}通过在 3DGS 隐式表面上构建图结构，利用 SelectionConv 表面卷积和预训练风格迁移网络直接对 .ply 或 .splat 文件进行风格化，CPU 上不到 2 分钟完成。FantasyStyle（2025.08）~\cite{2508_08136}是首个完全基于扩散模型蒸馏的 3DGS 风格迁移框架，通过多视角频率一致性机制缓解风格先验冲突。

\paragraph{DS-StyleGaussian 与 Style4D-Bench: 数据与基准}

风格化领域还有多种值得关注的新方向，以下从数据基准、特定风格和用户控制三个维度加以梳理。在\textbf{数据与基准}方面，DS-StyleGaussian（2026.05）~\cite{2605_30065}利用 2D 预训练解码器弥补 3D 风格迁移的数据稀缺问题，Style4D-Bench（2025.08）~\cite{2508_19243}则提出首个 4D 风格化基准，为时序一致性评估提供标准。

\paragraph{Thinking Like Van Gogh 等: 特定风格类型}

在\textbf{特定风格类型}方面，Thinking Like Van Gogh（2026.01）~\cite{2601_10075}将 2D 绘画流动场反向传播到 3D 高斯空间实现后印象派笔触风格，AI 驱动的 3D 环境风格化（2024.11）~\cite{2411_06067}将基本 3D 基元自动转化为风格化家具，GT2-GS（2025.05）~\cite{2505_15208}实现几何感知的纹理迁移。

\paragraph{Tune-Your-Style 与 DiffStyle3D: 用户控制}

在\textbf{用户控制}方面，Tune-Your-Style（2026.02）~\cite{2602_00618}支持用户灵活调整风格注入强度，DiffStyle3D（2026.01）~\cite{2601_19717}在扩散模型潜空间直接优化注意力特征。

\paragraph{TexGS-VolVis: 体积可视化风格化}

此外，TexGS-VolVis（2025.07）~\cite{2507_13586}将风格化扩展到体积可视化领域~\cite{2604_09639,2605_00498}。

\subsubsection{场景修复与物体移除}

本路线介绍在多视角一致前提下移除 3DGS 场景中的指定物体并修复被遮挡背景的方法。场景修复（Inpainting）和物体移除（Object Removal）是 3DGS 编辑中最实用的任务之一——用户标记想移除的物体后，系统需要多视角一致地恢复被遮挡的背景内容。这一任务的困难在于：不同视角看到的被遮挡区域不同，各自独立修复的结果可能互相矛盾；离散高斯表示下无法直接对修复区域进行梯度传播。

\paragraph{GScream: 交叉注意力特征传播修复}

GScream（2024.04）~\cite{2404_13679}是首个系统解决 3DGS 物体移除问题的方法。直觉上，它的工作方式类似于修复一面破损的墙——先用深度估计恢复墙壁的几何形状（把洞补平），再从周围完好的墙面"借"纹理涂到新补的区域上。其核心思路是从几何和纹理两个维度分别修复：先用单目深度引导改善几何，再通过交叉注意力（Cross-Attention）将可见区域的纹理传播到修复区域。

\textbf{动机与背景。}传统 NeRF 物体移除方法（如 SPIn-NeRF）依赖体渲染的连续表示，可以自然地通过多视角伪 RGB 引导进行梯度传播。然而 3DGS 的离散高斯表示打破了这一路径——每个高斯基元独立存在，修复区域的高斯与周围可见区域之间缺乏显式连接，无法直接从多视角修复图像中获取一致的梯度信号。此外，3DGS 的底层几何通常带有噪声，进一步加剧了物体移除后的几何空洞问题。

\textbf{深度引导的几何修复。}GScream 基于 Scaffold-GS 的锚点表示构建场景。为改善移除区域的几何质量，该方法引入单目深度估计（2024.03）~\cite{2403_12721}作为额外监督信号。具体地，对参考视角执行 2D 修复后，利用深度估计模型提取修复图像的深度图 $D_0$，同时提取其余视角图像的深度图 $D_i$。由于单目深度估计并非度量深度，GScream 设计了在线深度配准机制：通过最小二乘求解缩放参数 $w$ 和平移参数 $q$，将渲染深度 $\hat{D}_i$ 对齐到估计深度。为处理参考视角（已修复，包含完整深度）与其余视角（包含待移除物体，深度不可靠）之间的不一致性，GScream 提出加权深度损失：

\begin{equation}
\mathcal{L}_{\text{depth}} = \frac{1}{HW} \sum \hat{M}_i \, \| (w\hat{D}_i + q) - D_i \|,
\end{equation}

其中权重掩码 $\hat{M}_i$ 对参考视角的全图深度施加完整监督，对其余视角仅在背景区域施加监督。额外添加深度全变差损失 $\mathcal{L}_{\text{tv}}$ 以约束深度平滑性。

\textbf{交叉注意力特征传播。}仅靠深度监督只能解决几何问题，修复区域的纹理仍需从周围可见内容中获取。先前方法通常生成伪 RGB 引导图来更新场景表示，但这种策略对修复图像质量敏感，且忽略了修复区域与可见区域之间的内在关联。GScream 的关键洞察是：利用 3DGS 的显式表示特性，直接在 3D 空间中将可见区域的纹理信息传播到修复区域。具体实现分为两步：

\begin{enumerate}
\item \textbf{锚点采样}：对每个视角，采样同时覆盖掩码内外区域的 2D patch，将 3D 高斯锚点投影到当前视角，按投影位置将锚点分为修复区域组 $f_{\text{in}}$ 和周围可见区域组 $f_{\text{sur}}$。
\item \textbf{双向交叉注意力}：对两组锚点特征执行双向交叉注意力，实现信息的双向传播：

\begin{equation}
f'_{\text{in}} = \text{CrossAttn}(Q=f_{\text{in}}, K=f_{\text{sur}}, V=f_{\text{sur}}),
\end{equation}
\begin{equation}
f'_{\text{sur}} = \text{CrossAttn}(Q=f_{\text{sur}}, K=f_{\text{in}}, V=f_{\text{in}}).
\end{equation}
\end{enumerate}

更新后的特征被写回对应锚点，随后通过 Scaffold-GS 的标准流程（神经高斯生长和可微渲染）生成最终图像。梯度通过注意力机制回传到可见区域的锚点特征，使得修复区域中具有相似特征的锚点也能获得有效更新。这种设计使修复区域与可见区域在特征空间中保持一致，从而实现多视角连贯的纹理修复。

\paragraph{3DGIC: 掩码缩减与跨视角投影}

GScream 的交叉注意力特征传播虽然在特征空间建立了修复区域与可见区域的关联，但依赖 Scaffold-GS 的锚点结构，且传播效果受限于锚点采样的覆盖质量。3DGIC（2025.02）~\cite{2502_11801}从这一局限出发，其关键观察是：GScream 的特征传播在修复区域边界处容易产生模糊，因为注意力机制难以精确区分"真正需要修复的像素"和"从其他视角已可见的背景像素"。为此，3DGIC 放弃了特征空间的传播策略，转而在几何空间利用渲染深度做跨视角可见性分析。具体地，对参考视角的掩码区域 $M_1$，3DGIC 将其余视角的可见背景像素通过渲染深度 $D_{2:K}$ 投影到 3D 空间，再反投影回参考视角，提取落在 $M_1$ 内的像素坐标——这些像素在其他视角中已可见，无需修复。通过确定性几何投影而非特征匹配，3DGIC 将原始掩码 $M_1$ 精确缩减为 $M'_1$，仅保留真正被所有视角遮挡的区域。修复阶段在参考视角对渲染图像和深度图执行 2D 修复后，将修复结果投影为 3D 点云，再反投影到其余视角作为监督信号，通过 LPIPS 感知损失实现跨视角一致性。

\paragraph{SplatFill: 联合监督与选择性修复}

3DGIC 的掩码缩减策略有效减少了待修复区域面积，但其跨视角投影仍依赖单一参考视角的修复质量——当参考视角本身的修复结果存在误差时，误差会传播到其他视角。SplatFill（2025.09）~\cite{2509_07809}从两个维度解决这一问题。在监督维度，SplatFill 设计了联合深度-对象损失：软深度聚类损失（SDCL）将渲染深度按区间分箱并在箱内施加自监督聚类，避免了单目深度与度量深度之间的尺度配准问题（GScream 的在线配准机制对此较敏感）；裁剪聚焦深度损失（CFDL）对修复区域施加高斯加权聚焦监督，提升局部几何精度；对象感知对比损失（OACL）为每个高斯附加 16 维可学习分割特征，通过 SAM2 生成的分割掩码进行对比学习，使同一物体的特征聚集、不同物体的特征分离，实现对象级控制。在迭代维度，SplatFill 提出选择性引导修复（SGI）：每轮渲染所有训练视角，比较渲染深度与单目深度的误差图，选择累计误差最大的视角 $v^*$ 并在其内部定位几何不一致区域，仅对这些子区域重新执行 2D 修复并加入参考集。这种"定点清除"策略避免了对已一致区域的不必要扰动，使训练速度提升 24.5\%。

\paragraph{MVInpainter: 多视角扩散修复}

生成式方法利用扩散模型的强大修复能力，但需要解决多视角一致性问题。此前主流方案基于 SDS 优化 3D 表示——用 2D 扩散模型引导渲染图像向高概率密度区域移动。然而 SDS 在修复场景中暴露两个根本缺陷：一是优化过程的随机性导致每次生成结果不可控，修复任务恰恰需要精确恢复被遮挡区域的内容；二是生成的纹理常出现过平滑和高饱和，视觉质量难以令人满意。

MVInpainter（2024.09）~\cite{2409_16938}是首个用多视角扩散模型替代 SDS 进行修复的方法，其核心思路是：既然 SDS 只能提供"粗略的方向引导"，不如直接训练一个能从多个视角同时生成一致修复结果的扩散模型。整个流水线分为三个阶段：

\begin{enumerate}
\item \textbf{粗几何先验}：先用 SDS 在用户指定的 3D 边界框内生成粗糙几何作为形状参考（注意：此处 SDS 仅提供初始形状，不作为最终输出）。
\item \textbf{多视角扩散修复}：沿边界框中心的圆形轨迹均匀采样 $n$ 个视角，对每个视角渲染背景图 $I^{bg}$、2D 编辑掩码 $M$ 和来自粗几何的深度参考图 $D$，将三者送入 MVInpainter 生成视角一致的修复结果。
\item \textbf{掩码感知 3D 重建}：用修复结果微调原始 3D 高斯，同时约束未编辑视角保持不变。
\end{enumerate}

MVInpainter 的网络结构由两个分支组成。

\textbf{多视角扩散分支（MVD）}基于预训练的 SVD（Stable Video Diffusion）构建——SVD 本质是一个图像到视频模型，输入一张条件图像 $I^c$，经 VAE 编码后复制 $n$ 份形成条件潜变量序列，与带噪声的潜变量拼接后送入 UNet 去噪。然而，裸 SVD 生成的帧轨迹和内容不可预测，无法保证与指定相机轨迹一致。为此引入\textbf{ControlNet 条件注入分支}：将深度图 $D$、背景图 $I^{bg}$ 与掩码 $M$ 的逐元素乘积、以及 $M$ 本身沿通道维度拼接为 ControlNet Hints：

\begin{equation}
H_i = [I^{bg}_i \odot M_i,\; M_i,\; D_i], \quad i = 1, \ldots, n
\end{equation}

Hints 经可训练 CNN 下采样后，通过零初始化卷积层注入 SVD UNet 的编码器和中间块。零初始化保证训练初期 ControlNet 不干扰预训练权重，逐步学习注入几何先验、背景外观和编辑区域的控制信号。

掩码感知重建的关键在于：MVInpainter 最多生成 14 个视角的修复结果，稀疏视角无法为完整场景背景提供充分监督，导致新视角合成时出现浮点伪影。为此，MVInpainter 在微调时同时优化编辑视角和原始训练视角，掩码外区域强制保持不变：

\begin{equation}
\mathcal{L}_{rec} = \sum_{i \in \mathcal{C}^{edit}} \mathcal{L}_{gs}(I^r_i, I'_i) + \sum_{i \in \mathcal{C}^{gt}} \mathcal{L}_{gs}(I^r_i (1-M_i), I^{gt}_i (1-M_i))
\end{equation}

其中 $\mathcal{L}_{gs}$ 为标准 3DGS 重建损失（$L_1$ + D-SSIM），$\mathcal{C}^{edit}$ 为编辑视角集合，$\mathcal{C}^{gt}$ 为原始训练视角集合。该策略显著减少浮点伪影，在稀疏修复视角下仍保持高质量重建。

\paragraph{AuraFusion360 与 LightHarmony3D: 扩散修复扩展}

沿多视角扩散修复路线，后续方法将能力扩展到更复杂的场景类型。AuraFusion360（2025.02）~\cite{2502_05176}是首个针对 360 度无界场景的参考式 3D 修复方法，通过深度感知未见掩码生成、自适应引导深度扩散和 SDEdit 细节增强实现高质量对象移除。LightHarmony3D（2026.03）~\cite{2603_29209}专门解决物体插入中的光照一致性问题，通过生成式 HDR 环境图预测和 PBR 引导的阴影合成确保物理一致的阴影整合。

\paragraph{HOMER: 单应性传播加速修复}

上述方法在修复质量上不断突破，但计算效率仍是实际部署的瓶颈。

HOMER（2025.01）~\cite{2501_17636}将 3D 物体移除运行时间降至领先基线的五分之一。直觉上，它的工作方式类似于剪纸模板——只需在一个视角仔细剪好掩码并完成修复，后续视角通过几何变换自动复用这一结果，无需每个视角都从头开始。核心思路是用\emph{单应性传播}替代逐视角独立修复——用户只需在一个源视角标记目标区域，后续所有视角的掩码和修复结果都由几何传播自动生成。

\textbf{单应性掩码传播。}前述方法对每个视角独立运行扩散修复，计算开销随视角数线性增长。HOMER 的轻量 HoMM（Homography-based Mask Matching）模块仅处理相邻视角对：先用 LoFTR 计算密集关键点匹配，再估计单应性矩阵 $H$。给定源视角的二值掩码 $M_i$，目标视角的粗掩码通过逆单应性变换获得：

\begin{equation}
M^W_{i+1}(x, y) = M_i\!\bigl(H_{i,i+1}^{-1}[x,\, y,\, 1]^T\bigr)
\end{equation}

该变换保持掩码拓扑结构，同时根据视角间透视变化调整几何形状。匹配对数从 $N_v(N_v-1)$ 降至 $N_v-1$，大幅削减计算量。

\textbf{自适应锚圆调整。}稀疏视角或大视差下，单应性估计不可靠，经多步传播后误差累积显著。HOMER 引入锚圆策略将粗掩码精化为精确分割：以粗掩码 $M^W_i$ 的质心 $C_i$ 为圆心、半径 $r$ 为可调参数构造锚圆，圆上采样点作为 SAM 的点提示。$r$ 通过梯度下降优化，优化目标同时约束全局对齐和局部形状：

\begin{equation}
\mathcal{L} = \bigl(1 - \text{IoU}(M_i^k,\, M^{W,k}_i)\bigr) + \lambda\, \text{SCD}(M_i^k,\, M^{W,k}_i)
\end{equation}

其中 SCD 为形状上下文距离（Shape Context Distance），捕捉细粒度形状相似度；$\lambda$ 平衡两项权重。IoU 保证整体对齐，SCD 防止形状退化。

\textbf{修复传播与关键视角策略。}精化掩码获得后，HOMER 将源视角的修复结果沿掩码区域经 $H$ 扭曲到邻近视角，而非对每个视角重新修复。为避免扭曲误差跨视角累积，每 $n$ 个视角（默认 $n=10$）选取一个关键视角，用 LaMa 直接修复；其余视角仅接收传播结果。该策略将扩散模型调用次数从 $N_v$ 降至约 $N_v/n$，是运行时间压缩至五分之一的关键。整个流程兼容 NeRF 和 3DGS 两种辐射场后端。

\paragraph{VISTA 与 RePaintGS: 鲁棒性提升}

沿效率与鲁棒性方向，后续方法从不同角度提升修复质量。VISTA（2025.04）~\cite{2504_17815}从另一个角度提升鲁棒性：通过可见性不确定性引导和场景概念学习，测量 3D 点在不同视角中的可见性不确定性来引导多视角互补信息的利用。RePaintGS（2025.07）~\cite{2507_08434}通过参考视角引导和修复置信度评估，利用参考修复结果扭曲到其他视角作为伪真值。

\paragraph{ROODI: 遮挡物体提取}

在\textbf{遮挡处理与物体提取}方面，ROODI（2025.03）~\cite{2503_10256}结合 Wasserstein 距离修剪和扩散模型修复，首次有效解决 3DGS 场景中带遮挡的物体提取问题。

\paragraph{GOR-IS 与 D3DR: 物理一致性修复}

在\textbf{物理一致性}方面，GOR-IS（2026.05）~\cite{2605_00498}首创内在空间 3D 物体移除框架，通过 PBR 材质分解和自定义 3DGS 光线追踪器实现物理一致的全局光照处理（LPIPS 改善 13\%、PSNR 提升 2dB）；D3DR（2025.03）~\cite{2503_06740}利用扩散模型隐式的光照理解能力，通过 DDS 目标优化 3DGS 参数实现自然一致的物体插入和谐化（PSNR 提升 2.0dB）。

\paragraph{修复评估与质量保证}

在质量评估方面，Remove360（2025.08）~\cite{2508_11431}和语义残差测量（2025.03）~\cite{2503_17574}首次系统研究 3DGS 中物体移除后的语义残留问题，揭示几何移除与语义擦除之间的显著差距。高保真修复方法（2025.07）~\cite{2507_18023}通过自动掩码精炼和区域级不确定性引导优化实现多视角一致且照片级真实的修复。DeclutterNeRF（2025.04）~\cite{2504_04679}通过联合多视角可学习相机参数优化和遮挡退火正则化，无需生成先验即可移除复杂遮挡。

\begin{takeaway}
风格化与场景修复从艺术化渲染走向物理一致的物体移除和场景补全。GOR-IS 首创 PBR 材质分解和 3DGS 光线追踪器实现物理一致的全局光照处理，Remove360 揭示了几何移除与语义擦除之间的显著差距。核心挑战是在移除物体后恢复物理正确的光照和阴影。
\end{takeaway}

%=====================================================================

\subsection{近期方法脉络：物体移除与修复、Multi-view stylization of Gaussian scenes、2D-3D inpainting with Gaussian guidance等}

这一组近期工作延续 3DGS scene editing、3DGS、3DGS stylization，主要把改动集中在增量改进、编辑、扩展等环节。下面只保留相对既有路线的变化，便于和前文方法直接对照。

\paragraph{Benchmarking Single-Step Inpainting Methods for Multi-Object 3D Gaussian Splatting Scenes}
Benchmarking Single-Step Inpainting Methods for Multi-Object 3D Gaussian Splatting Scenes（2026.05）~\cite{2605_30987} 延续 3DGS scene editing 在“物体移除与修复”上的基本设定，主要改变是：多物体3DGS inpainting基准，对比2D修复器、从零重建和微调式3DGS修复策略。3DGS 的显式高斯表示和可微渲染能力保持为基础。
其设计动机是：构建多物体3DGS修复评测设置，系统比较单步inpainting流程中2D修复器和3DGS场景更新策略对跨视角一致性的影响。
具体而言，论文将2D inpainting输出接入3DGS场景创建和微调流程，比较先移除目标再修复、从头训练修复场景与对已有场景微调等策略。
适用边界是：主要是benchmark和经验分析；修复效果受2D inpainter质量和遮挡复杂度限制。

\paragraph{Capacity-Controlled Multi-View Stylization of 3D Gaussian Splatting}
Capacity-Controlled Multi-View Stylization of 3D Gaussian Splatting（2026.06）~\cite{2606_26754} 延续 3DGS 在“Multi-view stylization of Gaussian scenes”上的基本设定，主要改变是：基于半平衡最优传输的局部风格匹配与容量控制。3DGS 的显式高斯表示和可微渲染能力保持为基础。
其设计动机是：把局部风格匹配重写为 semi-balanced optimal transport，并加入可调列容量约束以稳定风格分配。
具体而言，针对每个渲染视图独立 2D feature matching 导致的视角不一致问题，在 3DGS 优化或编辑过程中约束风格特征的跨视角分配容量。
适用边界是：最优传输约束可能增加优化成本；极端风格或复杂遮挡场景的局部一致性仍需关注。

\paragraph{CoIn}
CoIn（2026.06）~\cite{2606_27584} 延续 3DGS 在“2D-3D inpainting with Gaussian guidance”上的基本设定，主要改变是：2D 扩散修复与 Reference Adaptive Gaussian Splatting 的多阶段一致性管线。3DGS 的显式高斯表示和可微渲染能力保持为基础。
其设计动机是：突破现有 GS 编辑方法依赖精确多视图 mask 且多用于物体移除的限制，支持对象插入等多样任务。
具体而言，先由 2D 扩散模型生成初始修复图像，再用 Reference Adaptive GS 和 Feature Attention 重建粗 3D 场景，并通过多阶段一致性优化减少跨视图冲突。
适用边界是：2D 扩散先验可能造成几何不一致或语义幻觉；多阶段优化成本和复杂遮挡区域仍是挑战。

\paragraph{Geometry-Aware Style Transfer in 3D Gaussian Splatting}
Geometry-Aware Style Transfer in 3D Gaussian Splatting（2026.06）~\cite{2606_24144} 延续 3DGS stylization 在“风格迁移与几何编辑”上的基本设定，主要改变是：Geometry-Aware Style Transfer：颜色与几何参数解耦交替优化，并用几何适配机制保持稳定。3DGS 的显式高斯表示和可微渲染能力保持为基础。
其设计动机是：将3DGS stylization从颜色迁移扩展到结构级几何适配，并通过解耦优化减少颜色和几何更新互相干扰。
具体而言，方法交替更新颜色和几何参数，在保持场景一致性的同时把风格中的结构属性迁移到高斯场。
适用边界是：几何风格迁移可能破坏真实度；风格和结构解耦程度需要调节。
\section{几何操控与其他方法}
\label{sec:09_3}

本节讨论 3DGS 的几何编辑能力，涵盖几何变形与操控（笼式变形、网格代理、拖拽式编辑等）以及动态场景编辑、跨模态应用等其他编辑方法。

\subsection{几何操控与扩展编辑}
\label{ss:09_3}

本节讨论两条技术路线：（1）几何变形与操控；（2）其他编辑方法。

\subsubsection{几何变形与操控}

本路线介绍通过笼式变形、网格代理和拖拽式操控等策略对 3DGS 场景进行几何变形的方法。3DGS 的离散高斯基元缺乏拓扑结构——没有面、边和连接关系。当需要将一个高斯场景"弯曲"或"拉伸"时，无法像网格模型那样直接移动顶点。几何变形方法需要解决两个核心问题：如何将连续变形场映射到离散高斯基元（包括位置、旋转和协方差的更新），以及如何在变形后保持纹理保真度和渲染质量。

\begin{table}[htbp]
\centering
\caption{3DGS 几何变形方法对比}
\label{tab:ch09_deformation}
\resizebox{\textwidth}{!}{
\begin{tabular}{lcccc}
\hline
\textbf{方法} & \textbf{变形策略} & \textbf{协方差更新} & \textbf{用户交互} & \textbf{兼容性} \\
\hline
GSDeformer & 笼式变形+代理点云 & 直接映射 & 笼顶点拖拽 & 通用 \\
CAGE-GS & 学习包围盒 & Jacobian $\Sigma'=J\Sigma J^\top$ & 多模态输入 & 通用 \\
Sketch-guided Cage & 笼式+Neural Jacobian & Neural Jacobian Field & 轮廓草图 & 静态场景 \\
Mani-GS & 三角网格代理 & 重心坐标绑定 & 网格变形 & 对物级 \\
D-MiSo & Triangle Soup & Core-Sub 层次 & 轨迹编辑 & 动态场景 \\
DragGaussian & 拖拽式 & 扩散引导 & 点对拖拽 & 通用 \\
DYG & 拖拽式+三平面 & Drag-SDS & 控制点拖拽 & 通用 \\
ARAP-GS & ARAP 能量 & 局部刚体约束 & 控制点对 & 通用 \\
SpLap & 拉普拉斯 & 泼溅图拉普拉斯 & 边界控制 & 通用 \\
FLOWING & 神经 ODE & 流场编码 & 关键帧 & 通用 \\
\hline
\end{tabular}
}
\end{table}

\paragraph{GSDeformer: 均值坐标笼式变形}

笼式变形（Cage-based Deformation）是经典的网格变形技术：构建包围目标的"笼"网格，用户拖拽笼顶点来驱动内部物体的变形。将其适配到 3DGS 面临一个根本挑战——3DGS 场景由离散的高斯基元构成，缺乏网格的拓扑连接关系，无法像顶点那样直接跟随笼变形。因此需要建立笼与离散高斯之间的映射关系。

GSDeformer（2024.05）~\cite{2405_15491}是首个笼式 3DGS 变形方法，其核心思路是将经典均值坐标（Mean Value Coordinates, MVC）适配到离散高斯基元。直觉上，均值坐标为笼内部的每个点分配一组权重，这些权重描述了该点的位置如何由笼顶点决定——笼顶点就像操控木偶的提线，移动任何一个顶点都会带动内部所有点按比例运动。

具体而言，GSDeformer 的变形管线分为三个阶段。第一阶段是\textbf{高斯→代理点云}转换：将每个高斯基元转换为显式椭球体，沿三个主轴与椭球面取六个交点作为代理点云 $AP = \{c, p_1^+, p_1^-, p_2^+, p_2^-, p_3^+\}$，使离散高斯获得可供变形操作的几何锚点。第二阶段是\textbf{MVC 变形}：对代理点云中的每个点，用源笼 $C_s$ 的笼顶点 $\{v_j\}$ 计算其笼坐标 $\{\phi_j\}$（即 $p = \sum_j \phi_j v_j$），然后直接替换为目标笼 $C_d$ 的笼顶点 $\{v'_j\}$，得到变形后位置 $p' = \sum_j \phi_j v'_j$。第三阶段是\textbf{变换反推与协方差更新}：由原始代理点 $AP_s$ 与变形后代理点 $AP_d$ 通过最小二乘拟合 $3 \times 3$ 变换矩阵 $T$，然后将高斯的均值更新为变形后中心，协方差更新为 $\Sigma_d = T \Sigma_s T^\top$。这一步至关重要——若不更新协方差，高斯的形状和朝向不会随变形调整，将产生严重的纹理拉伸伪影。

对于变形引起的极端弯曲，GSDeformer 引入\textbf{分裂过程}：当变形后高斯沿某轴的代理点弯曲角低于阈值时，将该高斯沿该轴一分为二，每个子高斯的轴长缩放 $k = \sqrt[3]{1/2} \approx 0.79$（近似保持总体积与形状），然后对子高斯重新执行 MVC 变形。这一策略解决了单个高斯无法表示弯曲形状的问题。

此外，GSDeformer 还设计了自动化笼生成管线：通过渲染多视角深度图 → T-SDF 重建 → 形态学闭运算 → 约束边折叠抽稀，从已训练的 3DGS 场景中自动构建低面数笼。其中关键创新是两阶段交替抽稀策略——先以约束 QEM 折叠边，再以 MVC 惩罚优化笼顶点位置使负权重趋零——确保笼的几何结构天然适配高质量变形。该方法兼容任何已训练的 vanilla 3DGS 或其变体，无需重训练即可实时编辑。

\paragraph{CAGE-GS: Jacobian 精确协方差更新}

沿笼式变形路线，CAGE-GS（2025.04）~\cite{2504_12800}从 GSDeformer 出发，针对其协方差更新的局限性提出改进。GSDeformer 通过代理点最小二乘拟合变换矩阵 $T$ 来更新协方差，但 $T$ 仅在代理点层面做全局近似，对局部变形场的描述精度有限。CAGE-GS 改用变形场的 Jacobian 矩阵来更新协方差：在高斯中心位置计算变形映射的 Jacobian $J = \partial p' / \partial p$（即变形后位置对原始位置的偏导），将协方差更新为 $\Sigma' = J \Sigma J^\top$。由于 Jacobian 直接刻画了每个高斯处的局部线性近似，能更精确地反映变形对高斯形状和朝向的影响，避免了 $T$ 拟合带来的近似误差。为降低逐高斯计算 Jacobian 的开销，CAGE-GS 先对采样子集计算 Jacobian，其余高斯通过 k-最近邻（kNN）继承最近采样点的 Jacobian。

此外，CAGE-GS 在笼的生成方式上也与 GSDeformer 不同。GSDeformer 通过自动化管线（多视角渲染 $\to$ T-SDF 重建 $\to$ 约束边折叠抽稀）从已有 3DGS 场景构建笼，而 CAGE-GS 借鉴 NeuralCage 的思路，训练神经网络从目标形状预测笼的偏移量，支持文本、图像、点云、网格和 3DGS 等多种输入模态。该方法移除了 NeuralCage 的预训练阶段，改为逐模型优化，使笼能更好地对齐目标形状。

\paragraph{Sketch-guided Cage: 草图驱动笼式变形}

沿笼式路线进一步推进交互性，Sketch-guided Cage（2024.11）~\cite{2411_12168}结合笼式变形与 Neural Jacobian Fields 实现细粒度控制，用户通过绘制轮廓草图即可指定变形方向，Neural Jacobian Fields 将草图转化为逐点的 Jacobian 矩阵场，实现比笼顶点拖拽更灵活的局部变形控制。ControlNet+SDS 确保变形结果的语义合理性，支持静态 3DGS 模型的动画化。

\paragraph{Mani-GS: 三角网格代理操控}

另一条路线利用三角网格作为变形的中间表示——网格拥有成熟的变形工具链，而 3DGS 可以通过绑定关系跟随网格变形。

Mani-GS（2024.05）~\cite{2405_17811}提出首个基于三角网格代理的 3DGS 操控框架。直觉上，它将高斯基元"粘"在三角形的局部坐标系上——就像把贴纸贴在三角形瓷砖上，移动或变形瓷砖时贴纸自动跟随，无需逐个重新定位。在此之前，SuGaR 等方法将高斯严格贴附于三角形表面（法线方向缩放趋零），渲染质量高度依赖网格精度——网格不准则高斯无法补偿缺失几何，网格虽精度高但仍有偏差时产生模糊伪影。更宽松的偏移绑定虽然允许高斯离开三角形，但偏移场不随变形迁移，操控后出现渲染扭曲。Mani-GS 的核心改进是\textbf{三角形感知绑定}（Triangle Shape Aware Binding）：不为每个高斯存储全局位置和旋转，而是在每个三角形 $f$ 的局部坐标系中优化局部参数。

具体而言，对于由顶点 $(v_1, v_2, v_3)$ 组成的三角形 $f$，Mani-GS 以第一条边方向为第一轴、三角形法向量 $\mathbf{n}^t$ 为第二轴、两者叉积为第三轴，构建局部坐标系 $R^t = [\mathbf{e}_1^t, \mathbf{n}^t, \mathbf{e}_3^t]$。高斯在三角形空间中的局部位置 $\mu^l$ 和局部旋转 $R^l$ 是可优化参数，全局属性通过坐标变换恢复：

\begin{equation*}
R = R^t R^l, \quad \mu = R^t \mu^l + c^t
\end{equation*}

其中 $c^t$ 为三角形中心的全局坐标。这样，当网格变形时三角形局部坐标系 $R^t$ 和中心 $c^t$ 自动更新，高斯的全局位置与旋转无需重新优化即可跟随变化。

此外，当三角形在变形中发生伸缩时，高斯的缩放也应按比例调整。Mani-GS 引入自适应向量 $\mathbf{e} = [l_1, \, 0.5(l_1+l_3), \, 0.5(l_2+l_3)]$（其中 $l_i$ 为三角形各边长度），将全局缩放修正为 $s = \mathbf{e} \odot s^l$，确保变形后渲染细节与表面形变一致。该绑定策略对网格精度有较高容忍度（即使使用 Marching Cube 提取的粗糙网格也能保持高保真渲染），并统一支持大变形、局部操控和软体模拟等多种操控类型。

\paragraph{Neural Surface Priors 与 D-MiSo: 网格代理增量}

沿网格代理路线，后续方法从不同角度扩展了能力。Neural Surface Priors（2024.11）~\cite{2411_18311}通过 PermutoSDF 获取高质量 SDF 作为几何先验约束高斯分布，设计网格到三角形汤代理的编辑传播机制，支持拓扑感知的编辑操作（如绑定、动画）。

D-MiSo（2024.05）~\cite{2405_14276}引入 Triangle Soup 表示动态高斯场景，通过 Core-Gaussians（大高斯，覆盖大空间范围）和 Sub-Gaussians（小高斯，附着于 Core-Gaussian 表面）的层次结构。每个 Sub-Gaussian 的位置由其所属 Core-Gaussian 的变换矩阵 $T$ 决定：$p'_{\text{sub}} = T \cdot p_{\text{sub}}$，因此修改 Core-Gaussian 即可同时改变所有关联 Sub-Gaussians，实现对动态场景中物体轨迹的独立编辑。

\paragraph{DragGaussian: 拖拽式编辑先驱}

笼式与网格代理方法均依赖额外的代理结构来传递变形，而拖拽式方法追求更直接的交互——用户指定控制点对，系统将高斯从源位置"拖"到目标位置，无需中间代理。

DragGaussian（2024.05）~\cite{2405_05800}是首个在 3D 高斯表示上实现拖拽式编辑的方法，利用扩散模型确保多视角一致性，支持开放词汇输入的交互式点对编辑。

\paragraph{DYG: 三平面拖拽编辑}

在 DragGaussian 的基础上，DYG（Drag Your Gaussian）（2025.01）~\cite{2501_18672}首次实现基于控制点拖拽的 3DGS 几何编辑。前述笼式与网格代理方法均依赖额外的代理结构来传递变形，DYG 的核心改进是去掉代理、直接在 3DGS 上拖拽，并通过隐式三平面表示弥补显式高斯的固有缺陷。其关键创新是\textbf{多分辨率三平面位置编码器}（Multi-resolution Triplane Positional Encoder）：将 3D 空间投影到三个正交平面 $(xy, xz, yz)$ 上，每个平面在多个分辨率下维护可学习特征图，任意位置 $p$ 的特征为三平面多尺度拼接
\begin{equation}
\mathbf{f}(p) = \phi\!\Big(\bigodot_{s} \bigodot_{c \in \{xy, xz, yz\}} \psi^s\!\big(H_c, \pi_c(x,y,z)\big)\Big),
\end{equation}

其中 $\pi_c$ 将点投影到第 $c$ 个平面，$\psi^s$ 在分辨率 $s$ 下进行双线性插值，$\odot$ 为 Hadamard 乘积，$\phi$ 为轻量级融合 MLP。这一设计解决了 3DGS 拖拽编辑中的核心困难——当用户将局部区域拖向新位置时，目标区域往往没有足够的高斯基元（稀疏），导致拖拽后的高斯无法稳定锚定；三平面作为连续隐式几何支架，为目标区域的空白位置提供合理的特征编码，使高斯可以在任意位置获得可靠的位移预测。

DYG 还设计了\textbf{区域感知位置解码器}：将高斯分为掩码内集合 $G_m$ 和掩码外集合 $G_{um}$，用两个 MLP 分别预测位移——掩码内高斯正常预测 $\Delta p = N_1(\mathbf{f})$，掩码外高斯通过残差修正 $\Delta p = N_1(\mathbf{f}) + N_2(\mathrm{sg}(\mathbf{f}))$（$\mathrm{sg}$ 为停止梯度算子）避免意外偏移，并辅以正则化损失 $L_\text{reg} = \sum_{g_i \in G_{um}} \|\Delta p_i\|$ 约束非编辑区域不动。整个拖拽过程分为两个阶段：第一阶段冻结所有高斯参数，仅优化三平面编码器和解码器，建立几何支架；第二阶段解冻高斯的颜色、不透明度、旋转和缩放属性并重新启用致密化操作，精化纹理细节。

在优化目标方面，DYG 提出\textbf{Drag-SDS} 损失函数，将用户拖拽意图编码为 2D 拖拽扩散模型（Lightning-Drag）的条件输入，通过 SDS 梯度同时优化高斯位置和三平面参数。Drag-SDS 由三项组成：潜空间损失 $L_{\text{lat}}$ 提供多视角一致的梯度引导，图像空间损失 $L_{\text{img}}$ 直接在像素域约束编辑质量，LoRA 适配损失 $L_{\text{lora}}$ 学习源分布的自适应估计以缓解 SDS 过平滑问题，总损失为 $L_{\text{Drag-SDS}} = \lambda_{\text{lat}} L_{\text{lat}} + \lambda_{\text{img}} L_{\text{img}} + \lambda_{\text{lora}} L_{\text{lora}}$。

\paragraph{ARAP-GS: 纯几何约束变形}

DYG 的拖拽流程依赖 Drag-SDS 损失通过扩散模型提供梯度引导，计算代价较高且 SDS 固有的过平滑可能损害细节。ARAP-GS（2025.04）~\cite{2504_12788}从另一个方向解决同一问题：放弃扩散引导，转而直接引入经典 ARAP（As-Rigid-As-Possible）能量作为几何先验来约束变形的局部刚性，从而在纯几何层面保证变形质量。经典 ARAP 能量要求每个顶点的局部变形近似刚体变换：

\begin{equation}
E_{\text{ARAP}} = \sum_i \sum_{j \in N(i)} w_{ij} \| (p'_i - p'_j) - R_i (p_i - p_j) \|^2,
\end{equation}

其中 $R_i$ 为最优旋转矩阵，$N(i)$ 为邻域索引集合。ARAP-GS 将高斯基元视为网格顶点，自然地将 ARAP 从网格推广到 3DGS——因为高斯中心的距离变化最小化等价于保持局部几何结构。

由于场景中高斯基元数量庞大（数十万级），直接对所有高斯求解 ARAP 系统在时间和内存上均不可行。ARAP-GS 为此设计了\textbf{子集采样-插值传播}策略：先对原始高斯集合进行随机采样得到代表性子集 $Q$，在 $Q$ 上通过 KNN 建立邻域关系 $N(i)$，在用户指定控制点对 $(h, h')$ 的约束下交替求解旋转矩阵 $R_i$（通过协方差矩阵的 SVD 分解）和变形位置 $p'_i$；然后将子集上的变形结果通过距离加权线性插值传播到全量高斯，包括位置插值 $p'_l = \sum_{i=1}^k w_{il} p'_i + p_l$ 和四元数旋转插值 $q'_l = \sum_{i=1}^k w_{il} q'_i \otimes q_l$（$\otimes$ 为四元数乘法），其中权重 $w_{il}$ 与到子集高斯的距离成反比。

ARAP 变形保证了几何刚性，但高斯的颜色和不透明度等属性在变形后可能产生伪影。ARAP-GS 在第二阶段引入\textbf{扩散先验外观精修}：利用现成的 StableSR 模型对渲染图像进行超分辨率增强，再以增强后的图像微调 3DGS 属性。为缓解多视角不一致问题，ARAP-GS 设计了三项策略：(1) 迭代数据集更新（IDU）——每 10 轮迭代仅替换一个子集视角的监督图像，而非一次性更新全部视角；(2) 掩码引导微调——根据高斯位移量生成编辑区域掩码 $M$，仅对编辑区域使用增强图像 $I_{\text{merge}} = M \odot I_{\text{sr}} + (1-M) \odot I_{\text{gt}}$，非编辑区域保持原始图像监督；(3) 深度损失——引入深度一致性约束以保持几何结构。这种两阶段设计（纯几何变形 + 扩散外观精修）将变形和外观优化解耦，避免了 DYG 中 SDS 梯度同时驱动几何和外观的耦合问题。

\paragraph{3DGS-Drag: 变形与 2D 编辑的融合}

3DGS-Drag（2026.01）~\cite{2601_07963}弥合变形方法和 2D 编辑方法之间的差距，利用 3DGS 的显式表示进行变形引导实现一致的几何修改，结合扩散引导进行内容校正和视觉质量增强。

\paragraph{SpLap: 泼溅图拉普拉斯变形}

除了上述三类主流变形策略，还有一些方法从不同角度重新定义变形问题。

\textbf{SpLap}（2025.11）~\cite{2511_19542}指出，现有 3DGS 变形方法要么依赖笼式/网格等外部代理（增加复杂度且受限于代理质量），要么将高斯基元直接当作点云处理（缺乏几何先验导致变形质量不佳）。SpLap 的核心洞察是：3DGS 自身的高斯核空间排布已蕴含表面几何信息，无需额外代理即可提取可靠的邻域结构。

为此，SpLap 提出\textbf{表面感知泼溅图}（Surface-Aware Splat Graph），直接在高斯基元上构建图结构。图的节点为每个高斯基元 $G_i$，边由高斯核的空间交集定义：首先为每个核计算\textbf{占用域}（occupancy region）$\Omega_i$——由激活阈值和最小贡献阈值约束的椭圆区域；然后通过法线方向偏移量 $\delta_{ij}$ 判定两个核是否相交，$\delta_{ij} < \epsilon$ 时建立连接。这种基于交集（而非单纯欧氏距离）的连通性定义，能有效过滤掉欧氏空间中距离近但沿表面实际很远的虚假邻居。

基于泼溅图，SpLap 的三阶段变形流水线为：(1) 在图上计算测地距离，结合 NonManifold 方法构建余切拉普拉斯算子 $L$；(2) 利用 $L$ 求解拉普拉斯变形系统（如 ARAP 或 BBW），得到变形映射 $\Phi$；(3) 为保持变形后的表面保真度，引入\textbf{核适配}阶段——为每个高斯核构造其占用域的最大内接三角形 $T_i$，通过距离反比加权将相邻高斯的位移传播到三角形顶点，再由 Steiner 外接椭圆唯一恢复适配后的核参数。该适配步骤确保变形后的高斯核仍紧贴表面流形，避免出现尖刺状伪影。

\paragraph{FLOWING: 神经 ODE 流场变形}

与 SpLap 在高斯基元上构建图结构不同，FLOWING（2025.10）~\cite{2510_10152}将变形重新构造为微分向量流的构建问题，通过基于流的神经架构——Neural ODE（NODE）建模连续变形轨迹、Neural Cost Flow（NCF）编码变形成本——在网络结构中自动保证连续性、可逆性和时间一致性，实现无需昂贵正则化的结构保持变形。

\paragraph{其他变形增量方法}

除了上述两类方法外，还有一些方法从不同角度扩展了 3DGS 变形的能力边界。在\textbf{输入模态}方面，单图像驱动的变形方法（2024.08）~\cite{2408_07540}引入位置导数（基于最优传输）捕获长距离运动，设计锚点 ARAP 正则化和粗到细优化策略处理遮挡高斯；视角依赖变形场（2025.04）~\cite{2504_05544}允许用户从不同视角对 3DGS 进行 2D 形变，通过平滑插值生成一致的 3D 变形场，支持类似图层编辑的堆叠变形。在\textbf{网格-高斯混合表示}方面，非刚性变形方法（2025.07）~\cite{2507_07000}将高斯核直接嵌入显式网格表面，利用网格固有的拓扑和几何先验引导编辑操作；SemMorph3D（2025.10）~\cite{2510_02034}通过 mesh 引导的 Gaussian 策略建立无监督语义对应，实现从图像到纹理完整的 3D 形状变形；ObjectMorpher（2026.03）~\cite{2603_28152}则将 2D 图像对象提升为可编辑 3DGS，使用 ARAP 约束确保物理合理的形状变化~\cite{2504_05544,2604_01551}。

\subsubsection{其他编辑方法}

本节介绍 3DGS 编辑中不属于上述主线方向但仍具重要价值的工作，包括动态与 4D 场景编辑和跨模态与特殊应用两个方向。

\paragraph{LoopGaussian 与 Gaussians-to-Life: 动态场景编辑先驱}

将静态 3DGS 编辑方法扩展到动态场景面临额外挑战：编辑需要在空间和时间两个维度上保持一致。

LoopGaussian（2024.04）~\cite{2404_08966}利用 3DGS 创建 3D Cinemagraph——通过 SuperGaussian 聚类和 Eulerian 运动场估计从多视角静态图像创建可循环的 3D 电影式动画。

沿动态编辑方向，Gaussians-to-Life（2024.11）~\cite{2411_19233}将能力从循环动画扩展到文本驱动动画，为任意预存在的 3DGS 场景添加文本驱动动画，利用视频扩散模型作为生成组件，不限于角色动画或单个物体。

\paragraph{CTRL-D: 单图驱动的动态场景编辑}

上述方法仅支持特定类型的动画生成，缺乏对动态场景进行精确编辑的能力。CTRL-D（2024.12）~\cite{2412_01792}提出首个利用单张编辑参考图像驱动整个动态场景编辑的框架。其核心思路是"用一张图教会扩散模型想要的编辑，再让它逐帧指导 3D 优化"。具体流程分三个阶段。（1）\textbf{单图编辑}：从输入的多视角图像中选取一张，用任意图像编辑工具（SDEdit、IP2P 或基于参考图的风格迁移）生成编辑版本。（2）\textbf{个性化 IP2P}：用编辑前后的图像对微调 InstructPix2Pix 的 U-Net，使其学会特定的编辑模式。由于只有一张配对图像，直接微调极易过拟合，因此 CTRL-D 引入两项对策：对源图和编辑图同时施加仿射变换（旋转、平移、剪切）做数据增强；借鉴 DreamBooth 的先验保留策略，用原始 IP2P 对随机帧生成伪编辑对作为正则项。微调损失为
\begin{equation}
  \mathcal{L}_{\text{finetune}} = \mathbb{E}_{\boldsymbol{\epsilon}, t}\!\bigl[\|\boldsymbol{\epsilon} - \boldsymbol{\epsilon}_\theta(z_t, t, I, C_T)\|^2\bigr] + \lambda\,\mathbb{E}_{\boldsymbol{\epsilon}, t}\!\bigl[\|\boldsymbol{\epsilon} - \boldsymbol{\epsilon}_\theta(z_t', t, I_d, C_T^0)\|^2\bigr],
\end{equation}

其中第一项在目标编辑对上训练，第二项保留模型的泛化能力，$\lambda$ 控制正则权重。（3）\textbf{两阶段动态场景优化}：第一阶段冻结变形场，仅用编辑关键帧优化规范空间的 3D 高斯参数并执行自适应致密化，填补编辑区域可能出现的几何空洞；第二阶段固定高斯数量，联合优化变形场与高斯属性。为保持帧间一致性，CTRL-D 维护一个编辑图像缓冲区——每次迭代随机选取一帧，用个性化 IP2P 生成编辑版本加入缓冲区，所有训练仅基于缓冲区内的图像。总损失融合重建损失与时序正则化：

\begin{equation}
  \mathcal{L} = (1 - \lambda_d)\,\mathcal{L}_1 + \lambda_d\,\mathcal{L}_{\text{D-SSIM}} + \lambda_t\,\mathcal{L}_{\text{temp}},
\end{equation}

其中 $\mathcal{L}_{\text{temp}}$ 最小化相邻帧间的差异，$\lambda_d = 0.2$，$\lambda_t = 0.001$。

\paragraph{4DGS-Craft 与 Dynamic-eDiTor: 4D 感知编辑}

在 CTRL-D 的单图驱动框架基础上，后续方法从不同维度提升 4D 编辑能力。4DGS-Craft（2025.10）~\cite{2510_01991}通过 4D 感知 IP2P 模型和 LLM 指令理解实现视角、时间和非编辑区域一致的 4D 场景编辑。Dynamic-eDiTor（2025.12）~\cite{2512_00677}利用多模态扩散变换器和时空子网格注意力机制实现无需训练的文本驱动 4D 场景编辑。

\paragraph{Catalyst4D: 3D 到 4D 的编辑迁移}

Catalyst4D（2026.03）~\cite{2603_12766}将高质量 3D 编辑转移到 4D 高斯场景。其核心是两个组件：\textbf{锚点运动引导}（AMG）在编辑前后的第一帧高斯上分别构建区域级锚点（通过 k-means 聚类），利用最优传输建立锚点间的语义对应 $\pi^* = \arg\min_\pi \sum_{i,j} \pi_{ij} c(a_i, a'_j)$，将第一帧的编辑变形传播到整个时序；\textbf{颜色不确定性外观精化}（CUAR）沿运动轨迹传播高置信度的外观信息，选择性修复因遮挡引起的颜色伪影。

\paragraph{CaricatureGS: 3D 人脸漫画化}

在动态编辑的特殊应用方面，CaricatureGS（2026.01）~\cite{2601_03319}通过曲率驱动的加权 Poisson 方程变形 FLAME 网格实现可控的 3D 人脸漫画化。

\paragraph{Gaussian Stitching: 示例建模与拼接}

编辑能力的扩展不止于变形和风格化，跨模态应用和特殊场景处理同样值得关注。

Gaussian Stitching（2024.08）~\cite{2408_15708}提出首个基于 3DGS 的真实感示例建模方法，通过基于采样的克隆优化实现多个 3DGS 模型的无缝拼接。

\paragraph{RaRa Clipper: 光滑连续裁剪}

沿跨模态应用方向，RaRa Clipper（2025.06）~\cite{2506_20202}将编辑能力扩展到裁剪操作，结合光栅化和光线追踪实现 3DGS 的平滑连续裁剪：对于完全在裁剪平面一侧的高斯使用高效光栅化，对于被裁剪平面截断的高斯切换到光线追踪，通过解析计算高斯函数与裁剪平面的交集确定截断区域，并用衰减函数 $\alpha' = \alpha \cdot \Phi(d/\sigma)$（$\Phi$ 为 sigmoid，$d$ 为距裁剪平面的有符号距离）调整边界处的不透明度，避免硬截断伪影。

\paragraph{DiT-Gaze: 视线重定向编辑}

在特殊应用方面，DiT-Gaze（2025.11）~\cite{2511_11231}首次在视线重定向中使用 Diffusion Transformer 渲染器，结合弱监督策略和正交约束损失实现高保真视线编辑。

\paragraph{可微分操作图与 Color3D: 建模与着色}

跨模态应用还有多个值得关注的方向。可微分操作图（2025.08）~\cite{2508_17645}将基本建模操作重构为可微分单元，首次解决 3D-AIGC 与设计师工作流的表示兼容性问题。Color3D（2025.10）~\cite{2510_10152}通过个性化着色器将单张关键视角图像的颜色信息传播到整个 3D 场景。物理引导分数蒸馏（2025.04）~\cite{2504_05296}为 3DGS 场景添加动态天气效果，通过物理仿真引导分数蒸馏联合精炼运动和外观。

\begin{takeaway}
几何操控与其他方法涵盖姿态编辑、组合编辑、网格提取后编辑和跨模态应用等方向。DiT-Gaze 将 Diffusion Transformer 引入视线重定向，可微分操作图解决了3D生成与设计师工作流的兼容性问题。这些工作共同推动3DGS编辑从单操作孤立处理走向多操作组合工作流。
\end{takeaway}

%=====================================================================

\subsection{近期方法脉络：对象转移与光照协调}

这一组近期工作延续 3DGS scene editing，主要把改动集中在增量改进等环节。下面只保留相对既有路线的变化，便于和前文方法直接对照。

\paragraph{Lighting-Consistent Object Transfer Across Radiance Fields}
Lighting-Consistent Object Transfer Across Radiance Fields（2026.06）~\cite{2606_22481} 延续 3DGS scene editing 在“对象转移与光照协调”上的基本设定，主要改变是：Lighting-consistent object transfer：用扩散模型协调源物体与目标场景的不一致光照合成图像。3DGS 的显式高斯表示和可微渲染能力保持为基础。
其设计动机是：提出面向radiance fields对象转移的光照协调模型，让提取对象粘贴到目标场景后更真实。
具体而言，方法训练扩散模型，把naive composite输入转化为光照一致输出，适用于VFX、室内设计和营销等场景。
适用边界是：扩散模型可能改变对象细节；3D几何和阴影一致性仍需要显式约束。

\section{近期研究（2026 年 8--11 月）}
\paragraph{DesignAgent3D：用交互规划与几何定位约束室内物体编辑}
2026.08~\cite{2608_21438}。DesignAgent3D 先以对话把含糊要求整理为明确编辑方案，再借 SfM 几何引导分割选定三维目标。编辑阶段用纹理参考与几何约束重绘选定视图，然后以替换视图和重训练的方式把改变写回场景表示，使编辑可从新视角观察。

论文同时测试 NeRF 与 3DGS 底座。对比编辑质量时，基线也收到已澄清的提示词，3DGS 方案在所测任务的 CLIP 方向相似度均值为 0.093。这个结果对应论文测试的室内场景与语义评分，比较本身不衡量基线澄清用户意图的能力。

\paragraph{CoGeo-GS：联合语义选取、深度补全与外观细化}
2026.08~\cite{2608_26656}。CoGeo-GS 将文本驱动的分割掩码蒸馏到高斯中，为多个目标分配语义身份。移除对象后，以对齐尺度的单目深度约束扩散补全，再用几何正则化外观细化缓和跨视图与边界不一致。

作者在 Mip-NeRF 360 和 SPIn-NeRF 上分别验证单对象和多对象移除，使用给定相机内参与 COLMAP 点云初始化，并与 InFusion、SPIn-NeRF、GaussianEditor 和 Gaussian Grouping 比较。被遮挡区域只能由条件补全，论文的视觉与重建指标针对该数据和相机设置。

\paragraph{LightBridge：将二维重光照提示传播到完整高斯场景}
2026.09~\cite{2609_02543}。LightBridge 先用学习到的光照变化潜空间在一步内预测重光照方向，再将二维提示传到可见高斯并沿三维表示传播到输入轨迹未覆盖的部分。模型固定高斯几何，只预测新的球谐颜色参数，避免为每个场景反复烘焙重光照视频。

论文分别评价原相机轨迹上的视频和最终高斯资产的新视角渲染。对只输出视频的基线，作者固定输入高斯几何并微调球谐系数 5000 步后比较。完全不可见的光源无法从输入掩码直接定位，静态合成室内训练的领域覆盖也限制了对新材质和灯具的泛化。

\paragraph{高斯调色板重参数化：在视图空间编辑颜色和亮度}
2026.09~\cite{2609_03897}。这项工作将预训练高斯的颜色重参数化为球谐调色板权重，渲染时混合出逐像素权重。高阶系数采用秩一分解，并以视图空间约束鼓励权重稀疏和颜色一致。交互时在固定高斯与权重下求解调色板与每色调曲线，用同一优化接收像素、调色板或曲线约束。

论文分别评估新视角重建保真度、编辑局部性和跨视图一致性，并与直接操作原语的像素编辑比较。结果可烘焙回标准 3DGS。由于共享调色板是控制界面，相同颜色的不同物体仍可能需要额外选择信息才能独立编辑。

\paragraph{三维编辑分解基准：用四个维度区分编辑成功与副作用}
2026.09~\cite{2609_14899}。该基准将指令遵循、目标局部性、跨视图一致性和非目标内容保持分别计分。八个代表性编辑器共享相机标定、原图、指令与目标标注，同类表示还共享初始重建，结果在受控留出视图上汇总。

作者在 240 组场景和编辑任务中发现，不同指标的领先方法不同。InterGSEdit 的指令保真度领先，而直接编辑代理的方案在三项保持导向指标领先。论文明确指出指标受视觉语言、分割和深度对应误差影响，静态、有空间目标的测试也不足以评估全局重光照或动态编辑。

\paragraph{φ-RIE：从静态高斯场景构造可操控的仿真资产}
2026.09~\cite{2609_26795}。φ-RIE 利用同一对象证据关联生成候选、资产对齐、原对象高斯移除与露出背景补全，将选定对象变成仿真中能移动的资产并保留场景其他部分。物理引擎输出的对象位姿继续驱动高斯外观渲染。

论文在 50 个 ScanNet++ 场景测试 1871 个对象请求。在固定保留与匹配数量的构造比较中，证据筛选与注册重试将作者报告的 F1 从 0.336 提高至 0.383，同时增加场景构造时间。严格验收后仍有匹配资产出现几何塌陷，说明通过筛选不保证几何正确或物理有效。

\paragraph{WINGS：利用三维生成先验修补删除对象后的高斯空缺}
2026.09~\cite{2609_37816}。WINGS 在用户给定的遮罩区域移除现有高斯，并以保留区域作为条件，借助预训练 TRELLIS 的结构潜变量模型生成空缺区域的原语。它处理的是无参考的条件生成，因此允许多种与现有场景相容的外观和几何。

论文在 COR-NeRF、Inpaint360 和 360-USID 上用 FID、KID、目标裁剪区域 FID、LPIPS 及用户偏好评价。一些复杂纹理在潜空间里偏向轴对齐，朝向不同的图案较难补全；论文将此列为结构旋转处理仍需改进的情形。

\paragraph{以可信关键帧锚定文本驱动的三维高斯编辑}
2026.10~\cite{2610_08179}。文本编辑若独立处理每个视图，较差视图的纹理可能将高斯场拉向不一致的结果。View Matters 先估计各视图的编辑重要性，偏向几何可见、语义明确且与指令相关的关键帧。编辑信息从关键帧传播到相邻视图，同时避免低可信分支的噪声反向传播，随后在高斯优化中延续该权重。

论文以 23 组场景和提示词的 CLIP 对齐及方向相似度比较编辑，并通过相邻视图与消融检查一致性。这类分数不能替代细部几何或边界真值。方法仍沿用冻结的二维编辑器，因此很大的几何改动和细粒度局部结构可能受到编辑器能力限制。

\paragraph{3DTexMOR：在共用纹理空间修补移除物体后的区域}
2026.10~\cite{2610_11198}。一次移除多个物体后，各视角暴露的遮挡区域并不一致。3DTexMOR 把跨视角的内在外观观测先聚合到同一个纹理空间，在该空间补全，再把结果投回各相机为高斯补全提供共用监督。附加的几何约束引导补全区域获得合理三维结构。

作者在合成的漫反射与反光场景中报告定量结果，在真实场景中查看视觉表现。方法采用共用平面近似，因此曲面、形变物体和形状不规则区域可能在反投影时产生错位，真实场景也没有可供量化对照的完整遮挡面真值。

\paragraph{RxGS：在可重光照高斯表示中单独编辑漫反射反照率}
2026.10~\cite{2610_25681}。RxGS 从带位姿多视图建成 BRDF 感知的高斯表示，只优化漫反射反照率以匹配参考图风格。粗糙度、金属度、法向和环境光保持原值，因而编辑后的场景仍能在新光照下重渲染。

作者用 NeRF-Synthetic 和 NeRF-LLFF 对比跨视图几何一致性，并加入 TensorIR 与 Mip-NeRF 360 检验重光照。几何一致性与重光照是分开的实验；原文指出参与前者的风格化基线不支持后者。该设计的适用范围受输入逆渲染分解和只编辑漫反射的约束。

\paragraph{VidEdit3D：沿相机轨迹传播锚点编辑}
2026.10~\cite{2610_35756}。VidEdit3D 将校准多视图排列为平滑的短视频轨迹，从用户提供的已编辑锚点出发，以预训练图生视频模型传播语义改变。最终区域选择约束把编辑回写到目标高斯，帮助减少对背景的改动。

论文表 1 在作者测试集上报告 CLIP-D 0.1275 和 MEt3R 0.1126，并在附录分析了遮挡、细结构、掩码不准和锚点内容不完整的失败例。未选高斯的参数可保持不变，但当目标形状改变遮挡关系时，它们的渲染可见性仍可能改变。

\paragraph{ReTouchGS：把场景级颜色操作直接写入高斯资产}
2026.11~\cite{2611_23054}。ReTouchGS 自动从预训练高斯及相机条件预测颜色操作，也支持将同一场景的原渲染与单张编辑参考配对后拟合操作。解析变换先修改原语颜色，再由混合感知的轻量细化网络修正结果，最终输出几何保持不变的标准高斯场景。

评测使用 37 个场景并以每场景 16 个离轨迹视图比较场景一致的颜色目标。度量反映对给定目标的保真度。免渲染描述的是共同的高斯参数更新阶段，参考迁移的前置颜色拟合仍需要原始渲染和编辑后的图像。

\paragraph{GS-VTON-MVT：以多视图联合编辑与几何条件开展高斯试衣}
2026.11~\cite{2611_25979}。GS-VTON-MVT 在给定已校准的多视图人物图像后重建个人的高斯场景。多视图 Transformer 先联合编辑一组视图来个性化扩散模型，随后用衣物图以及当前高斯场景渲染的深度与法向控制迭代编辑，并随相机姿态调节参考信息的影响。

论文用含 60 位受试者的 3D-VTONBench，每人选择 4 个训练视图进行多视图编辑，将其余视图留作新视角评价，并比较 FID、KID、基于 LPIPS 的一致性和用户偏好。这个实验协议依赖多视图输入与已建场景，不能据此推广至单张图像试衣。


\section*{本章小结}
\addcontentsline{toc}{section}{本章小结}

\begin{futurebox}{本章总结与展望}{ch09_summary}
本章综述了 3DGS 场景编辑与操控领域的 126 篇论文，涵盖以下核心方向：
\begin{itemize}
\item \textbf{文本引导编辑}：从 SDS 迭代优化到迭代数据集更新、渐进式定位和架构创新四条技术路线，代表方法包括 GaussianEditor、DreamCatalyst、ProEdit、VF-Editor 等
\item \textbf{交互式操控}：从 Point'n Move 的自提示分割到 GaussianBlock 的部件感知重建，再到 iSegMan 的免训练交互式操控
\item \textbf{风格化}：从优化式风格迁移到零样本前馈风格化，再到 CLIPGaussian 的统一多模态框架，以及 WaSt-3D 的分布匹配新范式
\item \textbf{场景修复}：从 GScream 的特征传播到 MVInpainter 的多视角扩散修复，再到 GOR-IS 的内在空间修复
\item \textbf{几何变形}：笼式（GSDeformer）、网格代理（Mani-GS）、拖拽式（DYG、ARAP-GS）和物理约束变形四条路线
\end{itemize}

当前该领域最突出的开放问题包括：
\begin{itemize}
\item \textbf{编辑精度与多视角一致性的根本矛盾}：扩散模型的大可行输出空间使精确编辑难以保持多视角一致，渐进式分解和蒸馏方法是当前最有效的缓解策略（详见第~\ref{ch:08}~章的生成方法讨论）
\item \textbf{复杂编辑的物理一致性}：变形后纹理拉伸、光照不一致、阴影缺失等问题仍缺乏统一解决方案，PBR 引导和内在空间处理是新兴方向
\item \textbf{4D 编辑的时空一致性}：动态场景编辑需要同时处理空间和时间维度的编辑传播，当前方法的效率和一致性均有提升空间
\item \textbf{编辑可逆性与版本管理}：大多数编辑方法不可逆，缺乏类似 Photoshop 的撤销/重做机制，ProGDF 的过程导向建模是初步尝试
\end{itemize}
\end{futurebox}
